%% 第26回編集委員会（物理第4巻）：組 電W7
%% p4-13.tex  電-13　長い電線を読む：二端子対回路と分布定数回路（★3〜5・3段）
%% 中身は計画書 tools/phys4-plan.md の B「電-13」。公式の持ち主は A（A81・A82）。中学・高校の公式は導かないので \phkoushiki は置かない（計画書 0 の1）。
%% 前の巻との約束（計画書 A-6）：P30（波-7 の 50 Ω・75 Ω と電気の線の反射：⑧・例題3・4。波-7 例題7 の r_V＝0.20・4 %、じぶんでとこうの 300 Ω→75 Ω の −0.60・36 %）、
%%   P40（波-14：ケーブルの中の電気の波：①⑥）、P43（ph3-back：インピーダンスと整合を電線の上の波で：⑧）。
%%   反射係数の書き方は波-7 ④（X-3）にそろえた：電圧（押すもの）の反射係数 r_V は波-7 の r_押 と同じ式、電流（流れるもの）の反射係数 r_I＝−r_V は波-7 の r と同じ式。
%%   記号 r_V・Z_先 は波-7 例題7 と同じ（波-7 の Z_線 がこの章の Z_0）。
%% 数学の本との約束（計画書 C-6-2）：7-3⑥・7-4:726（四端子定数の箱 (2 2；0.5 1)、AD−BC＝1）→ ③・例題7。7-11:366・9-11:1088（送電線の cosh・sinh、
%%   e^{Mℓ}、r＝1 Ω/km・g＝1 S/km・1 km → 1.543・1.175）→ ⑨・例題12。7-11 の z・y・ℓ は、この章の R'+jωL'・G'+jωC'・l（長さは第1巻の決まりで l）。
%% ほかの章の結果（導き直さない。計画書 B の節の番号で指す）：同軸の C'＝2πε/ln(b/a) は電-3 ⑩（例題7 と同じ寸法 0.45 mm・1.6 mm・εr 2.3 → 101 pF/m）、
%%   L'＝(μ0/2π)ln(b/a) は電-8 ⑥（自己誘導とインダクタンス。間の磁場は電-6 ⑤）、フェーザ・インピーダンスは電-10 ③〜⑤、三相は電-11 ⑦、変圧器は電-12 ⑤。
%%   波動方程式と c＝fλ は波-5 ②④、浅い水の波の2つの式（形をそろえた）は波-5 ⑤、ダランベールの解は 8-13⑨、cosh・sinh は 6-11⑤。
%% 数値はすべて Python で計算した：checks/物理第4巻/電-13/calc_p4_13.py（出力 calc_p4_13_out.txt）、
%%   ladder_p4_13.py（⑥の数値実験。本文のプログラムと同じ。出力 ladder_p4_13_out.txt、図のデータもここから）、kanji_p4_13.py（常用漢字の確かめ）。
%% 記号（この章）：1 m あたりの抵抗・インダクタンス・コンダクタンス・静電容量 R'・L'・G'・C'（Y-4）、特性インピーダンス Z_0（Y-4）、
%%   線の上の波の速さ c（W-3。真空の光速は c_0）、伝搬定数 γ＝α＋jβ（減衰定数 α・位相定数 β）、線の長さ l、位置 x、箱の長さ Δx、
%%   電圧・電流の瞬時値 v・i、フェーザ（複素数で書いた電圧・電流）V・I（電-10 と同じ）、電源の電圧 E（Y-3）、先につないだもののインピーダンス Z_先、
%%   電圧の反射係数 r_V・電流の反射係数 r_I、定在波比 SWR、四端子定数は行列 F＝(A B；C D) のまま扱う（Y-4：成分の C は静電容量ではないと③で断る。
%%   この章で静電容量は 1 m あたりの C' だけを使い、コンデンサー1個の C は使わない）、同軸の内半径 a・外半径 b、比誘電率 εr、
%%   短い線の抵抗・リアクタンス R・X、力率 cos φ、百分率インピーダンス %Z、定格電流 I_n、短絡電流 I_s。
%%   磁場は記号 B で書かない（四端子定数の B とぶつけない。「磁場」とことばで書く）。
%% 出典（第26回：書くときに確かめた。11月5日。WebSearch で、一次資料を引く資料を2つ以上照らし合わせた）
%%   ・地上デジタルテレビ放送の周波数 470〜710 MHz（UHF 13〜52 ch）：総務省「地上デジタルテレビジョン放送の一部周波数の変更」
%%       https://www.soumu.go.jp/menu_news/s-news/01ryutsu08_01000026.html と、総務省の資料（000105264.pdf）の周波数配置の図
%%   ・ポリエチレンの比誘電率 約2.3（2.25〜2.3）、RG-58（50 Ω）の速度係数 0.66：Farnell・Klaus Faber の RG-58 のデータシート（Velocity of propagation 66 %、
%%       solid PE）、everythingrf の表（PE 2.3・0.659）。発泡ポリエチレンは 1.5〜1.6 くらい（速度係数 0.78〜0.82）：例題2 では「1.5 とします」と書いた
%%   ・テレビ受信用の同軸ケーブル 75 Ω（JIS C 3502）・5C-FB は発泡ポリエチレン：帝国通信工業・光洋社（KTD）のカタログ。
%%       型番の読み方（5＝外部導体までの直径 約5 mm、C＝75 Ω・D＝50 Ω、F＝発泡ポリエチレン、B＝アルミはくテープと編組の2重）：かえるや（sds-kaeruya.com）の S-5C-FB の仕様書、weblio の用語解説
%%   ・LAN ケーブル（カテゴリー6）100 Ω：Helukabel・Farnell・EPS のデータシート（ISO/IEC 11801、TIA-568 を引く）。
%%       USB 約90 Ω・HDMI 約100 Ω（2本1組の線の差動の値）：Microchip の設計資料、TI の技術フォーラム、JLCPCB の設計資料（規格書そのものは会員向けで読めず）
%%   ・架空送電線の特性インピーダンス 300〜500 Ω、1相 約1 mH/km・約0.011 μF/km：台湾の公務員試験（輸配電）の問題の数（1.045 mH/km・0.01085 μF/km）、
%%       送電線のサージインピーダンスの解説（300〜500 Ω）。本文では「約300〜400 Ω」とし、例題には使わない
%%   ・6.6 kV 配電：電気学会論文誌の 6.6 kV 配電線のがいしの論文（J-STAGE ieejpes 139/7・132/9）、高圧は交流 600 V をこえ 7000 V 以下（電気設備技術基準の区分：
%%       パナソニックの用語解説で照合）。例題8 の線の抵抗・リアクタンスは「〜とします」
%%   ・%Z と短絡電流 I_s＝I_n×100/%Z：電験の解説（denken-ou.com）、viox の解説。例題9 の 30 kVA・5 % は「〜とします」
%%   ・トムソン（のちのケルビン男爵）"On the Theory of the Electric Telegraph"（Proc. R. Soc., 1855）：信号のおくれは長さの2乗に比例（抵抗と静電容量だけの線）。
%%       大西洋の電信ケーブル 1858 年（約3週間で使えなくなった）・1866 年に成功：atlantic-cable.com の Thomson の章、UNCW の講演資料
%%       （people.uncw.edu/hermanr/Documents/Talks/HeavisideThomsonTelegraphy.pdf）、バージニア大学の講義ノート（galileoandeinstein.phys.virginia.edu）
%%   ・ヘビサイド：1880年代（The Electrician の連載 1885〜1887）に、インダクタンスともれを入れた式を整え、ひずまない条件を示した。
%%       線にコイルを足す（装荷）は 1900 年ごろ実用（ピューピン・キャンベル）：EBSCO Research Starters、encyclopedia.com、ITM Web of Conferences 2019 の論文
%%   ・フェランチ効果：1890 年、ロンドンのデプトフォード発電所の 10 kV・約 11 km（7 マイル）の同軸の地下ケーブル。受電端の電圧が上がり、議論になった。
%%       フレミングが 1890〜91 年に調べた：ケンブリッジ大学ホイップル博物館の収蔵品の記録、UCL の文書館（フレミングの計算 1890〜1891）、Nature 44, 542 (1891)
%% 〔要確認〕なし（例題の配電線・変圧器・ケーブルの数は〔仮〕として本文で「〜とします」と書いた）
\makeatletter
% 抵抗（横）  \sxnR{(左はし)}{長さ}
\newcommand\sxnR[2]{\begin{scope}[shift={#1}]\filldraw[fill=white, draw=bkInk, line width=0.8pt] (0,-1.3) rectangle (#2,1.3);\end{scope}}
% 抵抗（たて）  \sxnRv{(上はし)}{長さ}
\newcommand\sxnRv[2]{\begin{scope}[shift={#1}]\filldraw[fill=white, draw=bkInk, line width=0.8pt] (-1.3,0) rectangle (1.3,-#2);\end{scope}}
% コイル（横、山4つ、長さ8）  \sxnL{(左はし)}
\newcommand\sxnL[1]{\begin{scope}[shift={#1}]\draw[bkInk, line width=0.8pt] (0,0) arc[start angle=180, end angle=0, radius=1] arc[start angle=180, end angle=0, radius=1] arc[start angle=180, end angle=0, radius=1] arc[start angle=180, end angle=0, radius=1];\end{scope}}
% コンデンサー（たて、まん中）  \sxnC{(まん中)}
\newcommand\sxnC[1]{\begin{scope}[shift={#1}]\draw[bkInk, line width=1.2pt] (-2.4,0.7) -- (2.4,0.7) (-2.4,-0.7) -- (2.4,-0.7);\end{scope}}
% 端子（白い丸）  \sxnT{(場所)}
\newcommand\sxnT[1]{\filldraw[fill=white, draw=bkInk, line width=0.7pt] #1 circle (0.9);}
% 同軸ケーブルの断面  \sxncoax{(中心)}{外の半径}
\newcommand\sxncoax[2]{\begin{scope}[shift={#1}]
  \fill[bkGray!55] (0,0) circle (#2+1.6);
  \fill[bkOrange!55!black] (0,0) circle (#2+0.5);
  \fill[bkYellow!25] (0,0) circle (#2);
  \fill[bkOrange!70!black] (0,0) circle (#2*0.28);\end{scope}}
% 画面の波形の枠  \sxnscreen{(左下)}{はば}{高さ}
\newcommand\sxnscreen[3]{\filldraw[fill=bkNavy!8, draw=bkInk!70, line width=0.6pt, rounded corners=0.8mm] #1 rectangle ++(#2,#3);}
\makeatother

\phchapter{13}{長い電線を読む：\\二端子対回路と分布定数回路}{テレビのアンテナ線に、なぜふつうの電線を使わない？}{\bklevelrange{3}{5}}{長い電線を読む：二端子対回路と分布定数回路}
\bkmichishirube{6-11,7-3,7-4,7-11,8-13,9-11,波-5,波-7,電-3,電-8,電-10}{電-15,電-16}

\begin{goal}
\begin{itemize}[label={\color{bkTheme}●}, leftmargin=1.6\zw, itemsep=0pt, topsep=1pt]
\item \bklevel{3}\enspace ケーブルの上を電気の波が進み、先やつなぎ目ではね返ることを、速さ・往復の時間・波長で説明できる。
\item \bklevel{3}\enspace 特性インピーダンスの違いから、電圧の反射係数と、はね返る電力の割合を求められる。
\item \bklevel{4}\enspace 四端子定数の行列で箱をつなぎ、短い配電線の電圧降下と、変圧器の短絡電流を見積もれる。
\item \bklevel{5}\enspace 小さな箱のくさりから電信方程式を導き、波の速さ $1/\sqrt{L'C'}$ と特性インピーダンス $\sqrt{L'/C'}$ を求められる。
\item \bklevel{5}\enspace 損失のある線・長い送電線・4分の1波長の線を、$\cosh$・$\sinh$ と四端子定数で扱える。
\end{itemize}
\end{goal}

\begin{hatena}
\noindent\begin{minipage}[c]{0.55\linewidth}\raggedright
屋根のアンテナからテレビまでは、黒くて丸い、少し太いケーブルでつないであります。電灯のコードのような、2本の線を並べたふつうの電線は使いません。

2本の線があれば、電気の信号は通るはずです。なぜ、わざわざ特別なケーブルを使うのでしょう。
\end{minipage}\hfill
\begin{minipage}[c]{0.42\linewidth}\centering
\begin{tikzpicture}[x=0.8mm, y=0.8mm]
  % 家とアンテナ
  \fill[bkOrange!15] (0,0) rectangle (26,18);
  \draw[bkInk!70, line width=0.6pt] (0,0) rectangle (26,18);
  \fill[bkRed!35] (-2,18) -- (13,28) -- (28,18) -- cycle;
  \draw[bkInk!70, line width=0.6pt] (-2,18) -- (13,28) -- (28,18);
  \draw[bkInk, line width=0.8pt] (20,22) -- (20,38);
  \foreach \yy in {30,33,36} {\draw[bkInk, line width=0.8pt] (16,\yy) -- (24,\yy);}
  % ケーブル
  \draw[bkInk, line width=1.6pt, rounded corners=2mm] (20,24) -- (31,24) -- (31,6) -- (24,6);
  % テレビ
  \fill[bkGray!40] (6,2) rectangle (22,13);
  \fill[bkNavy!15] (7.2,3.2) rectangle (20.8,11.8);
  \bkperson{(12,3.6)}{7.4}{bkBlue}
  \begin{scope}[opacity=0.35]\bkperson{(15.6,3.6)}{7.4}{bkBlue}\end{scope}
  \node[font=\sffamily\footnotesize, anchor=west, align=left] at (33,15) {同軸\\ケーブル};
  \node[font=\sffamily\footnotesize, anchor=north] at (14,-0.6) {ゴースト（二重の像）};
\end{tikzpicture}
\end{minipage}
\end{hatena}

\noindent 第1〜3部では、電池や家のコンセントにつないだ回路を考えました。そこでは、導線の上のどこでも電流は同じで、スイッチを入れると、回路全体の電圧と電流が「同時に」決まる、として計算してきました（電-4・電-5・電-10）。ところが、電気の信号も、実は有限の速さで電線を伝わっています。線が長くなったり、信号の変わり方が速くなったりすると、「同時」とは言えなくなり、電線の上を\textbf{電気の波}が進み、はね返るようすが見えてきます。この章では、そのような長い電線を、波として読む方法と、行列（四端子定数）で読む方法を学びます。物理シリーズの『波動』（波-7）で学んだ、ひもや音の「インピーダンスと反射」が、そのまま電気の線でも使えます。

\begin{yattemiyou}[なわの上のパルスと、家のケーブルの文字]
\textbf{用意するもの}：長いなわとびのなわ（\qty{3}{m}くらい）、家にあるテレビのケーブル・LAN ケーブル（見るだけ）\par\vspace{2pt}
\noindent\begin{minipage}[c]{0.64\linewidth}
\begin{enumerate}[label={\bkcircled[bkGreen]{\arabic*}}, leftmargin=2.4\zw, itemsep=1pt, topsep=0pt]
\item なわの先を、友だちにしっかり持ってもらう（動かない先）。手元を1回すばやく横にふって、山を送る。はね返る山は、同じ側か、反対側か。
\item 友だちに、なわの先を軽く持って、自由に動かしてもらう（自由に動く先）。同じように山を送る。
\item テレビのケーブルや LAN ケーブルの外側に印刷された文字を読む。「75Ω」「5C-FB」「CAT6」などの文字があるか。
\end{enumerate}
\end{minipage}\hfill
\begin{minipage}[c]{0.33\linewidth}\centering
\begin{tikzpicture}[x=0.85mm, y=0.85mm]
  \draw[bkInk, line width=1.0pt] (0,10) -- (40,10);
  \draw[bkBlue!80!black, line width=1.3pt] (6,10) -- ++(4,6) -- ++(4,-6);
  \draw[bkBlue!80!black, line width=0.7pt, -{Stealth[length=1.8mm]}] (6,19) -- (16,19);
  \fill[bkInk] (40,10) circle (1.1);
  \node[font=\sffamily\footnotesize, anchor=north] at (34,8) {先};
  \fill[bkGray!55] (8,-2) rectangle (40,2);
  \node[font=\sffamily\footnotesize, text=bkInk] at (24,0) {5C-FB 75Ω};
\end{tikzpicture}
\end{minipage}\par\vspace{2pt}
\noindent\textbf{注意}：ケーブルは見るだけにします。コンセントにつないだ機器のケーブルを抜いたり、切ったり、中を開けたりしてはいけません。
\end{yattemiyou}

\noindent やってみようの答え：先を動かないように持つと、山は\textbf{反対側}（谷）になってはね返ります。先が自由に動くと、山は\textbf{同じ側}のままはね返ります（波-6 ①・波-7 ⑤の固定端と自由端）。テレビのケーブルには「75Ω」や「5C-FB」、LAN ケーブルには「CAT6」などが印刷されていることが多いでしょう。75 Ω は、ケーブルの長さによらない、ケーブルの\textbf{特性インピーダンス}という値です。この章で、その意味と、なわのはね返りと同じことが電線でも起きることを学びます。

\begin{omoidasou}
\bkrecall{波-7}{インピーダンス $Z$＝押すもの÷流れるもの（弦は力÷速さ、音は圧力÷速さ）。境目では、速さ（流れるもの）で数えた反射係数 $r=\dfrac{Z_1-Z_2}{Z_1+Z_2}$、押すもので数えた反射係数 $r_{\text{押}}=\dfrac{Z_2-Z_1}{Z_1+Z_2}=-r$。はね返るエネルギーの割合は $r^2$。}
\bkrecall{波-5}{波動方程式 $\dfrac{\partial^2y}{\partial t^2}=c^2\dfrac{\partial^2y}{\partial x^2}$ の答えは、形を保って速さ $c$ で進む波。$c=f\lambda$。}
\bkrecall{電-10}{交流の電圧・電流を複素数（フェーザ）$V$・$I$ で書くと、コイルは $V=j\omega L\,I$、コンデンサーは $I=j\omega C\,V$。インピーダンス $Z=V/I$。}
\bkrecall{7-3}{入口2本・出口2本の線の箱を、出口の $(V_2,\ I_2)$ を入れると入口の $(V_1,\ I_1)$ が出る2×2 の行列で書く。箱をつなぐと行列のかけ算（7-3⑥）。}
\bkrecall{6-11}{$\cosh x=\dfrac{e^x+e^{-x}}{2}$、$\sinh x=\dfrac{e^x-e^{-x}}{2}$。$y''=a^2y$ の答えは $\cosh ax$ と $\sinh ax$ の組み合わせ（6-11⑤）。}
\end{omoidasou}

%% ======================== 物理の道具 ========================
\bksection{doguu}{物理の道具}

\phstage{me}{電線の上を進む電気の波を見る}{3}

\bksubsection{1}{電線の上を、電気の波が進む}
\label{ph:電-13:sec1}
長さ\qty{10}{m}のテレビのケーブルの入口に、ほんの一瞬だけ電圧をかける装置（パルス発生器）をつなぎます。入口とケーブルの先の電圧を、とても速い変化まで見える電圧計（\textbf{オシロスコープ}）で見てみましょう。



\noindent 入口で出したパルスは、すぐには先に着きません。約\qty{50}{ns}かかって先に着き、先ではね返って、また約\qty{50}{ns}かけて入口に戻ってきます（図\ref{ph:電-13:fig-tdr}）。往復\qty{20}{m}を約\qty{100}{ns}なので、速さは
\[\frac{\qty{20}{m}}{\qty{100e-9}{s}}=\qty{2.0e8}{m/s}\]
で、真空の光の速さ $c_0=\qty{3.00e8}{m/s}$ の約3分の2です。\textbf{電気の信号は、電線の上を、光の速さに近い有限の速さで進む}のです。この章では、線の上の波の速さを $c$、真空の光の速さを $c_0$ と書きます（物理シリーズの『波動』と同じく、波の速さを $c$ と書くため）。

\begin{bkfigure}
\begin{tikzpicture}[x=1mm, y=1mm]
  % 装置とケーブル
  \filldraw[fill=bkBlue!12, draw=bkInk, line width=0.7pt, rounded corners=1mm] (0,30) rectangle (20,42);
  \node[font=\sffamily\footnotesize, align=center] at (10,36) {パルス\\発生器};
  \draw[bkInk, line width=2.0pt] (20,36) -- (104,36);
  \draw[bkGray!70, line width=1.0pt] (20,36) -- (104,36);
  \fill[bkInk] (104,36) circle (1.2);
  \node[font=\sffamily\footnotesize, anchor=north] at (62,34.2) {長さ\qty{10}{m}の同軸ケーブル（\qty{50}{\ohm}）};
  \node[font=\sffamily\footnotesize, anchor=west, align=left] at (106,36) {先};
  \draw[bkBlue!80!black, line width=1.2pt] (32,36) -- (34,41) -- (36,36);
  \draw[bkBlue!80!black, line width=0.7pt, -{Stealth[length=1.8mm]}] (38,40) -- (48,40);
  \node[font=\sffamily\footnotesize, anchor=west] at (49,40.5) {約\qty{50}{ns}で先へ};
  % 3つの画面
  \foreach \k/\lab/\sgn in {0/{先を開ける\\（何もつながない）}/1, 1/{先をショート\\（2本をつなぐ）}/-1, 2/{先に\qty{50}{\ohm}の\\抵抗をつなぐ}/0} {
    \begin{scope}[shift={(\k*43,0)}]
      \sxnscreen{(0,0)}{38}{22}
      \draw[bkGray!60, line width=0.3pt] (2,8) -- (36,8);
      \draw[bkRed!80!black, line width=1.0pt] (2,8) -- (6,8) -- (6,15) -- (8,15) -- (8,8) -- (26,8) -- (26,{8+7*\sgn}) -- (28,{8+7*\sgn}) -- (28,8) -- (36,8);
      \node[font=\sffamily\footnotesize, anchor=north, align=center] at (19,-1) {\lab};
    \end{scope}}
  \node[font=\sffamily\footnotesize, anchor=south] at (7,16) {出した};
  \node[font=\sffamily\footnotesize, anchor=south, fill=bkNavy!8, inner sep=1pt] at (27,15.6) {戻った};
  \draw[bkInk, line width=0.4pt, {Stealth[length=1.5mm]}-{Stealth[length=1.5mm]}] (7,3) -- (27,3);
  \node[font=\sffamily\footnotesize, anchor=south, fill=bkNavy!8, inner sep=1pt] at (17,3.2) {約\qty{100}{ns}};
\end{tikzpicture}
\bkfigcaption{入口の電圧をオシロスコープで見た模式図（横は時刻）。出したパルスは、約\qty{100}{ns}（1 ns は10億分の1秒）後に、先ではね返って戻ってくる。先を開けると同じ向き、ショートさせると反対向きで戻り、\qty{50}{\ohm}の抵抗をつなぐと戻らない。}\label{ph:電-13:fig-tdr}
\end{bkfigure}

\begin{phyomi}[電子がその速さで走るのではない]
速さ\qty{2.0e8}{m/s}で進むのは、電子そのものではありません。導線の中の電子が動く速さは、1秒に1 mm にもならないほどゆっくりです（電-4 ②）。進むのは、「押された」という知らせ、つまり電圧と電流の\textbf{形}です。ちょうど、なわの1つ1つの部分はその場でゆれるだけなのに、山の形が先へ進むのと同じです（波-5 ①）。電気の知らせを運んでいるのは、2本の線のあいだの空間にできる電場と磁場で、その速さは、線のあいだを満たす物（ケーブルではプラスチック）で決まります（⑦）。
\end{phyomi}

\phboxhead{はね返り方は、先のようすで変わる}
図\ref{ph:電-13:fig-tdr}の3つの画面を、やってみようのなわと比べましょう。
\begin{itemize}[leftmargin=1.2\zw, itemsep=1pt, topsep=1pt, label={\color{bkTheme}●}]
\item \textbf{先を開ける}：電流が先から出ていけないので、電荷が先にたまって、電圧のパルスは\textbf{同じ向き}のまま戻ります。
\item \textbf{先をショート}：2本の線をつなぐと、先の電圧は0にしかなれないので、電圧のパルスは\textbf{反対向き}になって戻ります（なわの先を動かないように持ったときと同じ）。
\item \textbf{先に\qty{50}{\ohm}の抵抗}：ケーブルと同じ値の抵抗をつなぐと、パルスは抵抗に吸いこまれて熱になり、\textbf{まったく戻りません}。
\end{itemize}
3つめが、ケーブルに書かれた「\qty{50}{\ohm}」の意味です。このケーブルの上を進む波は、電圧と電流の比がいつも\qty{50}{\ohm}になっています。先に\qty{50}{\ohm}の抵抗をつなぐと、抵抗は、波から見て「ケーブルがまだずっと続いている」のと区別がつかないので、はね返しようがないのです。この比を、ケーブルの\textbf{特性インピーダンス}といい、$Z_0$ と書きます。\bkidx{とくせいいんぴーだんす}{特性インピーダンス}波-7 ③の「押すもの÷流れるもの」で、電気では、押すものが電圧、流れるものが電流です。

\phboxhead{はね返る割合は、インピーダンスの比で決まる}
先に、$Z_0$ と違うインピーダンス $Z_{\text{先}}$ の機器をつなぐと、波の一部がはね返ります。はね返る波の電圧と来た波の電圧の比（\textbf{電圧の反射係数}）$r_V$ は、波-7 ④で、力や音圧（押すもの）で数えた反射係数 $r_{\text{押}}$ と同じ式で計算できます（式の導き方は⑧）。
\pheq{r_V=\frac{Z_{\text{先}}-Z_0}{Z_{\text{先}}+Z_0}\qquad\text{はね返る電力の割合}=r_V^{\,2}}{単位なし}
先を開けると（$Z_{\text{先}}\to\infty$）$r_V=+1$、ショートさせると（$Z_{\text{先}}=0$）$r_V=-1$、$Z_{\text{先}}=Z_0$ なら $r_V=0$ で、図\ref{ph:電-13:fig-tdr}の3つの画面のとおりです。\qty{50}{\ohm}のケーブルに\qty{75}{\ohm}の機器をつなぐと、$r_V=25/125=0.20$ で、電力の4\,\% がはね返ります（波-7 例題7、\bkreidairef{3}）。

\phboxhead{「短い線」と「長い線」}
では、家の電灯の配線では、なぜこのような波を考えなくてよかったのでしょう。比べるのは、線の長さと、信号の\textbf{波長}です（$c=f\lambda$：波-5 ②）。家のコンセントの電気は、1秒に50回か60回、向きが変わる交流です。線の上の速さを $c_0$ に近いとして、\qty{50}{Hz}の波長は
\[\lambda=\frac{c}{f}\approx\frac{\qty{3.00e8}{m/s}}{\qty{50}{Hz}}=\qty{6.0e6}{m}=\qty{6000}{km}\]
です（\bkreidairef{6}）。家の配線の長さ（数十 m）は、波長の10万分の1ほどです。波の山1つが、日本列島ほどの長さなので、家の中の線の上では、どこでも電圧はほとんど同じ時刻に同じように変わります。このような線を「\textbf{短い線}」とよび、これまでのように、線の上のどこでも電流は同じ、として計算できます。部品が1か所ずつにまとまった回路と見てよいので、\textbf{集中定数回路}ともいいます。\bkidx{しゅうちゅうていすうかいろ}{集中定数回路}

テレビの電波は\qtyrange{470}{710}{MHz}です。ケーブルの中の速さを\qty{2.4e8}{m/s}くらいとすると、波長は\qtyrange{0.34}{0.51}{m}で、屋根からテレビまでの\qty{10}{m}のケーブルの上には、波の山が20〜30個も並びます。このような線が「\textbf{長い線}」で、線の上の電圧は場所ごとに違い、波として扱わなければなりません。線の長さが波長のおよそ10分の1をこえたら、長い線として考える、というのが工学のおおよその目安です。

\begin{phbaisuru}{周波数 $f$}
波長 $\lambda=c/f$ は半分になります。同じ長さの線でも、周波数を上げていくと、短い線から長い線に変わります。
\par\smallskip\noindent\begin{tikzpicture}
\begin{axis}[bkplot, axis lines=left, width=104mm, height=44mm, xmode=log, ymode=log, xmin=10, xmax=1e10, ymin=1e-2, ymax=1e8, xtick={1e2,1e4,1e6,1e8,1e10}, xticklabels={100 Hz,10 kHz,1 MHz,100 MHz,10 GHz}, ytick={1e-2,1,1e2,1e4,1e6,1e8}, yticklabels={1 cm,1 m,100 m,10 km,1000 km,$10^5$ km}, xlabel={周波数 $f$}, ylabel={波長 $\lambda$}]
  \addplot[bkcurve, domain=10:1e10, samples=20] {2.0e8/x};
  \fill[bkRed] (axis cs:50,4e6) circle (1.8pt) (axis cs:600e6,0.333) circle (1.8pt) (axis cs:1e9,0.2) circle (1.8pt);
  \node[font=\sffamily\footnotesize, anchor=south west, fill=white, inner sep=1pt] at (axis cs:70,5e6) {家の電気 50 Hz};
  \node[font=\sffamily\footnotesize, anchor=north east, fill=white, inner sep=1pt] at (axis cs:2e8,0.8) {テレビ・スマホ（数百 MHz〜数 GHz）};
  \draw[bkGray, line width=0.4pt, dashed] (axis cs:10,10) -- (axis cs:1e10,10);
  \node[font=\sffamily\footnotesize, anchor=south west, text=bkGray, fill=white, inner sep=1pt] at (axis cs:15,11) {家の配線の長さ（約 10 m）};
\end{axis}
\end{tikzpicture}
\par\noindent（速さ\qty{2.0e8}{m/s}の線の上の波長。たても横も10倍ごとに同じ長さの目もり。）
\end{phbaisuru}

\bksubsection{2}{同軸ケーブル・平行2線・基板の線}
\label{ph:電-13:sec2}
電気の波を運ぶ線は、どれも2本の導体でできています。電流が行く道と帰る道です。2本のあいだの空間に電場と磁場ができ、それが信号を運びます。いろいろな線の形を見てみましょう（図\ref{ph:電-13:fig-katachi}）。

\begin{bkfigure}
\begin{tikzpicture}[x=1mm, y=1mm]
  % 同軸
  \sxncoax{(14,14)}{9}
  \foreach \a in {0,45,...,315} {\draw[bkRed!80!black, line width=0.7pt, -{Stealth[length=1.4mm]}] (\a:3.2) ++(14,14) -- ++(\a:5.0);}
  \draw[bkBlue!80!black, line width=0.6pt, dashed] (14,14) circle (6.6);
  \node[font=\sffamily\footnotesize, anchor=north, align=center] at (14,0) {同軸\\ケーブル\\50 Ω・75 Ω};
  % 平行2線
  \begin{scope}[shift={(44,14)}]
    \fill[bkGray!25, rounded corners=2mm] (-12,-2.5) rectangle (12,2.5);
    \fill[bkOrange!70!black] (-8,0) circle (1.5) (8,0) circle (1.5);
    \draw[bkRed!80!black, line width=0.6pt, -{Stealth[length=1.4mm]}] (-6.2,0) -- (6.2,0);
    \draw[bkRed!80!black, line width=0.6pt, -{Stealth[length=1.4mm]}] (-6.8,1.2) .. controls (-3,6) and (3,6) .. (6.8,1.2);
    \draw[bkRed!80!black, line width=0.6pt, -{Stealth[length=1.4mm]}] (-6.8,-1.2) .. controls (-3,-6) and (3,-6) .. (6.8,-1.2);
  \end{scope}
  \node[font=\sffamily\footnotesize, anchor=north, align=center] at (44,0) {平行2線\\（フィーダー）\\300 Ω};
  % より線
  \begin{scope}[shift={(74,14)}]
    \draw[bkBlue!80!black, line width=1.4pt, domain=-11:11, samples=60] plot (\x, {2*sin(30*\x)});
    \draw[bkOrange!80!black, line width=1.4pt, domain=-11:11, samples=60] plot (\x, {-2*sin(30*\x)});
  \end{scope}
  \node[font=\sffamily\footnotesize, anchor=north, align=center] at (74,0) {より線の組\\（LAN）\\約 100 Ω};
  % 基板の線
  \begin{scope}[shift={(106,14)}]
    \fill[bkGreen!25] (-13,-4) rectangle (13,1);
    \fill[bkOrange!70!black] (-13,-5.2) rectangle (13,-4);
    \fill[bkOrange!70!black] (-2.5,1) rectangle (2.5,2.2);
    \draw[bkRed!80!black, line width=0.6pt, -{Stealth[length=1.4mm]}] (0,1) -- (0,-3.8);
    \draw[bkRed!80!black, line width=0.6pt, -{Stealth[length=1.4mm]}] (-2.2,1) -- (-5,-3.8);
    \draw[bkRed!80!black, line width=0.6pt, -{Stealth[length=1.4mm]}] (2.2,1) -- (5,-3.8);
  \end{scope}
  \node[font=\sffamily\footnotesize, anchor=north, align=center] at (106,0) {基板の線\\（マイクロストリップ）\\約 50 Ω};
\end{tikzpicture}
\bkfigcaption{電気の波を運ぶ、いろいろな線の断面（同軸は輪切り、より線は横から見た図）。実線の矢印は電場（電気力線）、同軸の破線の円は磁場（磁力線）の向きの1本。同軸ケーブルでは、電場と磁場が、外の導体の内側にとじこめられる。}\label{ph:電-13:fig-katachi}
\end{bkfigure}

\phboxhead{同軸ケーブル}
中心に1本の導線（内導体）があり、その周りを、プラスチック（ポリエチレンなど）の絶縁体、網や薄い金属の筒（外導体）、外側のカバーが、同じ中心で取りまいています。\textbf{同軸}は「軸が同じ」という意味です。\bkidx{どうじくけーぶる}{同軸ケーブル}信号の電流は、内導体を行き、外導体の内側の面を帰ってきます。電場は内導体から外導体へ放射状に、磁場は内導体を囲む円の形にでき、どちらも\textbf{外導体の内側にとじこめられます}。だから、
\begin{itemize}[leftmargin=1.2\zw, itemsep=1pt, topsep=1pt, label={\color{bkTheme}●}]
\item 信号が、電波になって外へもれにくい。
\item 外から来る雑音の電波が、中に入りにくい（外導体が盾になる。電-3 ③の静電遮へいと似たはたらき）。
\item 太さの比と絶縁体が決まれば、特性インピーダンスが決まり、近くに金属や人の手があっても変わらない。
\end{itemize}
テレビのケーブルは\qty{75}{\ohm}、無線機や測定器のケーブルは\qty{50}{\ohm}のものが多く使われます（波-7 ⑧）。

\begin{phmodel}[電線を、波を運ぶ線と見る]
\phignore{導線の抵抗と、絶縁体の中のわずかな電流のもれ（⑨で入れる）。線の曲がり。線のはしの形。}
\phassume{2本の導体の形と、あいだの絶縁体が、線にそってどこも同じ（一様な線）。}
\phlimit{線が途中でつぶれたり、曲げすぎたりした所（そこで $Z_0$ が変わり、はね返る）。周波数がとても高く、断面の大きさが波長と同じくらいになるとき（電-15 ⑩の導波管の世界）。}
\end{phmodel}

\phboxhead{平行2線・より線・基板の線}
2本の線を決まった間かくで平らなひもの中に並べた\textbf{平行2線}（フィーダー）は、昔のテレビのアンテナ線に使われました（約\qty{300}{\ohm}）。電場と磁場が外の空間に広がっているので、近くの物の影響や雑音を受けやすく、今はほとんど同軸ケーブルに変わりました。LAN ケーブルの中には、2本の線をらせんのように\textbf{より合わせた組}が4組入っています。より合わせると、外の雑音が2本の線に同じように入り、2本の電圧の差だけを受け取る機器では打ち消し合います（特性インピーダンスは約\qty{100}{\ohm}）。スマホやパソコンの電子回路の\textbf{基板}の上の細い銅の線も、裏の銅の面と組になって、特性インピーダンスが約\qty{50}{\ohm}になるように、線のはばと板の厚さが設計されています。USB のケーブルは約\qty{90}{\ohm}、HDMI のケーブルは約\qty{100}{\ohm}です。

\begin{phmijika}[ケーブルの数字を読む]
テレビのケーブルの「5C-FB」は、日本の決まった書き方で、「5」は外導体の内側の直径（絶縁体の太さ）がおよそ\qty{5}{mm}、「C」は\qty{75}{\ohm}、「F」は絶縁体が泡を入れたポリエチレン（発泡ポリエチレン）、「B」は外導体がアルミはくをはったテープと網の2重であることを表します。\qty{50}{\ohm}のケーブルは「5D-2V」のように「D」で書きます。同じ\qty{75}{\ohm}でも、太いケーブルほど、信号が弱まりにくくなります。
\end{phmijika}



\Needspace{12\baselineskip}\phstage{te}{箱の行列で、線と変圧器を計算する}{4}

\bksubsection{3}{四端子定数：箱を行列で表し、つなぐ}
\label{ph:電-13:sec3}
長い線の話にもどる前に、線や変圧器を「箱」として計算する道具を用意します。ケーブル、送電線、変圧器、雑音を取りのぞくフィルターは、どれも、電気が入る2本の線（入口）と、出ていく2本の線（出口）をもつ箱と見ることができます。端子の対（2本1組）が2つあるので、このような箱を\textbf{二端子対回路}（にたんしついかいろ）といいます。\bkidx{にたんしついかいろ}{二端子対回路}

\begin{bkfigure}
\begin{tikzpicture}[x=1mm, y=1mm]
  \draw[bkInk, line width=0.9pt] (0,10) -- (24,10) (0,-10) -- (24,-10) (56,10) -- (80,10) (56,-10) -- (80,-10);
  \filldraw[fill=bkTheme!10, draw=bkDeep, line width=1.1pt, rounded corners=2mm] (24,-14) rectangle (56,14);
  \node[font=\sffamily\small, text=bkDeep] at (40,0) {箱 $\mathbf F$};
  \sxnT{(0,10)}\sxnT{(0,-10)}\sxnT{(80,10)}\sxnT{(80,-10)}
  \draw[bkOrange!85!black, line width=1.2pt, -{Stealth[length=3.5pt]}] (6,13) -- (16,13);
  \node[font=\sffamily\footnotesize, anchor=south] at (11,13.5) {$I_1$};
  \draw[bkOrange!85!black, line width=1.2pt, -{Stealth[length=3.5pt]}] (64,13) -- (74,13);
  \node[font=\sffamily\footnotesize, anchor=south] at (69,13.5) {$I_2$（出ていく向き）};
  \draw[bkBlue!80!black, line width=0.8pt, {Stealth[length=3pt]}-] (-4,8) -- (-4,-8);
  \node[font=\sffamily\footnotesize, anchor=east] at (-5,0) {$V_1$};
  \draw[bkBlue!80!black, line width=0.8pt, {Stealth[length=3pt]}-] (84,8) -- (84,-8);
  \node[font=\sffamily\footnotesize, anchor=west] at (85,0) {$V_2$};
  \node[font=\sffamily\footnotesize, anchor=north] at (8,-11.5) {入口};
  \node[font=\sffamily\footnotesize, anchor=north] at (72,-11.5) {出口};
  \node[anchor=west, font=\small] at (94,0) {$\begin{pmatrix}V_1\\I_1\end{pmatrix}=\mathbf F\begin{pmatrix}V_2\\I_2\end{pmatrix}$};
\end{tikzpicture}
\bkfigcaption{二端子対回路。出口の電圧・電流 $(V_2,\ I_2)$ を入れると、入口の $(V_1,\ I_1)$ が出てくる2×2 の行列 $\mathbf F$ で表す（7-3⑥と同じ向き）。電流 $I_2$ は、出口から次の箱へ出ていく向きを正にとる。}\label{ph:電-13:fig-hako}
\end{bkfigure}

\noindent 交流の電圧・電流を複素数（フェーザ）$V$・$I$ で書きます（電-10 ③）。数学の 7-3⑥では、出口の組 $(V_2,\ I_2)$ を入れると入口の組 $(V_1,\ I_1)$ が出てくる行列で、箱を表しました。電気の本では、この行列の4つの成分を $A$・$B$・$C$・$D$ と書いて、\textbf{四端子定数}とよびます。\bkidx{よんたんしていすう}{四端子定数}
\pheq{\begin{pmatrix}V_1\\I_1\end{pmatrix}=\begin{pmatrix}A&B\\C&D\end{pmatrix}\begin{pmatrix}V_2\\I_2\end{pmatrix}}{$A$・$D$ は単位なし、$B$ は Ω、$C$ は S}
この行列を $\mathbf F$（F 行列）と書きます。成分の単位は、$A$ が「電圧÷電圧」、$D$ が「電流÷電流」で単位なし、$B$ が「電圧÷電流」で Ω、$C$ が「電流÷電圧」で S（ジーメンス）です。

\noindent\textbf{文字の注意}：四端子定数の $C$ は、コンデンサーの静電容量ではありません。この章では、静電容量は「1 m あたりの静電容量」$C'$ だけを使い、四端子定数は、なるべく行列 $\mathbf F$ のまま扱います。成分の $B$ も、磁場（電-6）とは関係ありません。

\phboxhead{2つの基本の箱}
上の線のとちゅうにインピーダンス $Z$ の部品を入れた箱（直列の箱）と、2本の線のあいだにアドミタンス $Y=1/Z$ の部品をわたした箱（並列の箱）は、7-3⑥の抵抗の箱と同じ考えで、
\pheq{\text{直列の箱}\ \begin{pmatrix}1&Z\\0&1\end{pmatrix}\qquad\text{並列の箱}\ \begin{pmatrix}1&0\\Y&1\end{pmatrix}}{Ω と S}
です。直列の箱では、電流は通りぬけるだけ（$I_1=I_2$）で、入口の電圧は部品の電圧 $ZI_2$ だけ高くなります（電-4 ⑦の電圧則）。並列の箱では、電圧は同じ（$V_1=V_2$）で、入口の電流は、部品を流れる $YV_2$ と出口へ行く $I_2$ の和です（電流則）。$Z$ と $Y$ は、抵抗なら $R$ と $1/R$、コイルなら $j\omega L$ と $1/(j\omega L)$ のように、どんな部品でもかまいません。

\phboxhead{箱をつなぐと、行列のかけ算：縦続接続}
前の箱の出口を、次の箱の入口につなぐことを\textbf{縦続接続}（じゅうぞくせつぞく）といいます。\bkidx{じゅうぞくせつぞく}{縦続接続}前の箱の出口の組は、次の箱の入口の組そのものなので、
\[\begin{pmatrix}V_1\\I_1\end{pmatrix}=\mathbf F_1\begin{pmatrix}V_2\\I_2\end{pmatrix}=\mathbf F_1\mathbf F_2\begin{pmatrix}V_3\\I_3\end{pmatrix}\]
で、つないだ箱全体の行列は、\textbf{入口側から順に並べた行列の積} $\mathbf F_1\mathbf F_2$ です。何個つないでも、かけ算を繰り返すだけです。7-3⑥の直列\qty{2}{\ohm}のあとに並列\qty{2}{\ohm}（$Y=\qty{0.5}{S}$）をつないだ箱では、行列は1行め $(2,\ 2)$、2行め $(0.5,\ 1)$ でした（\bkreidairef{7}）。

\phboxhead{先に何かをつなぐと：入口から見たインピーダンス}
箱の出口に、インピーダンス $Z_{\text{先}}$ の機器をつなぐと、$V_2=Z_{\text{先}}I_2$ です。入口から見たインピーダンス $Z_{\text{入}}=V_1/I_1$ は、
\pheq{Z_{\text{入}}=\frac{AV_2+BI_2}{CV_2+DI_2}=\frac{AZ_{\text{先}}+B}{CZ_{\text{先}}+D}}{Ω}
\label{ph:電-13:zin}です（分子と分母を $I_2$ でわった）。箱をいくつつないでも、全体の行列の4つの成分がわかれば、入口から見たインピーダンスが1つの式で計算できます。

\phboxhead{$AD-BC=1$}
直列の箱も並列の箱も、行列式（7-4②）は $1\times1-Z\times0=1$、$1\times1-0\times Y=1$ です。積の行列式は行列式の積（7-4⑤）なので、基本の箱をいくつつないでも、
\pheq{AD-BC=1}{単位なし（Ω×S）}
が成り立ちます（7-4⑦）。\textbf{抵抗・コイル・コンデンサー・変圧器だけでできた箱では、いつも $AD-BC=1$} です。計算した4つの数がこれを満たしていなければ、どこかで計算を間違えた、と気づくことができます。

\begin{seikaku}
$AD-BC=1$ は、「入口に電源をつないで、ショートさせた出口の電流をはかっても、出口に同じ電源をつないで、ショートさせた入口の電流をはかっても、同じになる」という性質（相反性）の表れです。トランジスタのように、電源から電気をもらって信号を大きくする部品の入った箱では、成り立ちません。また、この章では電流 $I_2$ を「出口から出ていく向き」にとりました。$I_2$ を「出口から箱へ入る向き」にとる本では、$B$ と $D$ の符号が変わり、$AD-BC=-1$ の形になります（7-3⑥のせいかくに言うと）。
\end{seikaku}

\bksubsection{4}{短い送電線・配電線のモデル}
\label{ph:電-13:sec4}
変電所から工場や家の近くまで電気を運ぶ\textbf{配電線}は、長さが数 km〜十数 km です。\qty{50}{Hz}の波長\qty{6000}{km}に比べてずっと短いので、①の「短い線」です。短い線では、線の抵抗 $R$ と、線の自己インダクタンスによるリアクタンス $X=\omega L$（電-10 ④）が、線全体に広がっているのを、1か所にまとめて「直列の箱」1つで表せます。2本の線のあいだの静電容量による電流は小さいので、無視します。

\begin{bkfigure}
\begin{tikzpicture}[x=1mm, y=1mm]
  % 回路
  \draw[bkInk, line width=0.8pt] (0,14) -- (8,14) (16,14) -- (20,14) (28,14) -- (40,14) (0,0) -- (40,0);
  \sxnR{(8,14)}{8}\sxnL{(20,14)}
  \node[font=\sffamily\footnotesize, anchor=south] at (12,15.6) {$R$};
  \node[font=\sffamily\footnotesize, anchor=south] at (24,16.2) {$jX$};
  \sxnT{(0,14)}\sxnT{(0,0)}
  \filldraw[fill=bkGreen!15, draw=bkInk, line width=0.7pt, rounded corners=1mm] (40,-2) rectangle (50,16);
  \node[font=\sffamily\footnotesize, align=center] at (45,7) {負\\荷};
  \node[font=\sffamily\footnotesize, anchor=east] at (-1.5,7) {$V_{\mathrm s}$};
  \node[font=\sffamily\footnotesize, anchor=west] at (51,7) {$V_{\mathrm r}$};
  \draw[bkOrange!85!black, line width=1.1pt, -{Stealth[length=3pt]}] (30,17) -- (37,17);
  \node[font=\sffamily\footnotesize, anchor=south] at (34,17.4) {$I$};
  \node[font=\sffamily\footnotesize, anchor=north, align=center] at (24,-2.5) {送電端から受電端へ};
  % フェーザの図
  \begin{scope}[shift={(62,6)}]
    \draw[bkBlue!80!black, line width=1.3pt, -{Stealth[length=2.6mm]}] (0,0) -- (40,0);
    \node[font=\sffamily\footnotesize, anchor=north] at (24,-0.6) {$V_{\mathrm r}$};
    \draw[bkOrange!85!black, line width=1.1pt, -{Stealth[length=2.4mm]}] (0,0) -- (16,-12);
    \node[font=\sffamily\footnotesize, anchor=north east] at (15,-11.5) {$I$};
    \draw[bkInk, line width=0.4pt] (8,0) arc[start angle=0, end angle=-36.87, radius=8];
    \node[font=\sffamily\footnotesize, anchor=west] at (8.4,-2.8) {$\phi$};
    \draw[bkRed!80!black, line width=1.1pt, -{Stealth[length=2.2mm]}] (40,0) -- (51.2,-8.4);
    \node[font=\sffamily\footnotesize, anchor=north east] at (46,-5) {$RI$};
    \draw[bkGreen!50!black, line width=1.1pt, -{Stealth[length=2.2mm]}] (51.2,-8.4) -- (62,6);
    \node[font=\sffamily\footnotesize, anchor=west] at (58.2,-3) {$jXI$};
    \draw[bkInk, line width=1.3pt, -{Stealth[length=2.6mm]}] (0,0) -- (62,6);
    \node[font=\sffamily\footnotesize, anchor=south east] at (30,3.2) {$V_{\mathrm s}$};
    \draw[bkInk, line width=0.4pt, {Stealth[length=1.4mm]}-{Stealth[length=1.4mm]}] (40,11) -- (62,11);
    \draw[bkGray, line width=0.4pt, dashed] (40,0.8) -- (40,12) (62,7) -- (62,12);
    \node[font=\sffamily\footnotesize, anchor=south] at (51,11.4) {$\Delta V$（ほぼ）};
  \end{scope}
\end{tikzpicture}
\bkfigcaption{短い線のモデル（左）と、電圧のフェーザの図（右。$RI$・$jXI$ は見やすいように大きくかいた）。電流 $I$ は、受電端の電圧 $V_{\mathrm r}$ より角 $\phi$ だけおくれている。$V_{\mathrm s}$ は、$V_{\mathrm r}$ に、$I$ と同じ向きの $RI$ と、$I$ より90°進んだ $jXI$ をたしたもの。}\label{ph:電-13:fig-mijikai}
\end{bkfigure}

\noindent 送電端（変電所の側）の電圧を $V_{\mathrm s}$、受電端（負荷の側）の電圧を $V_{\mathrm r}$、線を流れる電流を $I$ とします。直列の箱の行列から、
\[V_{\mathrm s}=V_{\mathrm r}+(R+jX)I\]
です。工場の機械（モーター）のように、電流が電圧より角 $\phi$ だけおくれる負荷では（電-11 ③の力率 $\cos\phi$）、図\ref{ph:電-13:fig-mijikai}のように、$RI$ は電流と同じ向き、$jXI$ は電流より90°進んだ向きです（$j$ をかけると90°回る：電-10 ③）。$V_{\mathrm r}$ の向きにそろえて見ると、$RI$ は $RI\cos\phi$、$jXI$ は $XI\sin\phi$ だけ、$V_{\mathrm r}$ を長くします。$V_{\mathrm r}$ に直角な向きのずれは、矢印の長さをほとんど変えないので、
\pheq{\Delta V=|V_{\mathrm s}|-|V_{\mathrm r}|\approx I(R\cos\phi+X\sin\phi)}{V ＝ A×Ω}
\label{ph:電-13:koka}です。これを\textbf{電圧降下}といいます。\bkidx{でんあつこうか}{電圧降下}

\begin{phyomi}[$\Delta V\approx I(R\cos\phi+X\sin\phi)$]
「線の電圧の下がりには、抵抗は力率 $\cos\phi$ の分だけ、リアクタンスは $\sin\phi$ の分だけ効く」。力率が1（電流と電圧が同じ向き。電熱器）なら $\Delta V=RI$ で、直流と同じです。電流がおくれるほど、$X$ の効きが大きくなります。送電線や配電線は、$X$ が $R$ と同じくらいか大きいので、力率が悪いと電圧がよけいに下がります。
\end{phyomi}

\begin{phbaisuru}{電流 $I$}
電圧降下も2倍です。夏の夕方など、たくさんの家や工場が電気を使うと、配電線の電流が増えて、変電所から遠い家ほど電圧が下がります。家の電圧が\qty{101}{V}の上下\qty{6}{V}の中に入るように（電気事業法施行規則 第38条）、変電所や電柱の上の機器で電圧を調整しています。
\end{phbaisuru}

\begin{phmodel}[短い線]
\phignore{2本の線のあいだの静電容量に流れる電流。線の上の波（往復の時間のずれ）。}
\phassume{線の長さが、交流の波長よりずっと短い（50 Hz なら数十 km まで）。三相の線は、1相分（1本の線と大地のあいだ）で考える（電-11 ⑦）。}
\phlimit{数百 km の長い送電線や、静電容量の大きい地中ケーブル（⑩）。}
\end{phmodel}

\bksubsection{5}{百分率インピーダンスと短絡電流}
\label{ph:電-13:sec5}
送電線や変圧器には、それぞれの\textbf{定格}（安全に使い続けられる電圧 $V_{\mathrm n}$・電流 $I_{\mathrm n}$）があります。大きさの違う機器を比べるとき、電力の工学では、インピーダンスを「定格の電流を流したとき、その機器の中で下がる電圧が、定格の電圧の何\,\% か」で表します。これを\textbf{百分率インピーダンス}（\%インピーダンス）といいます。\bkidx{ひゃくぶんりついんぴーだんす}{百分率インピーダンス}
\pheq{\%Z=\frac{ZI_{\mathrm n}}{V_{\mathrm n}}\times100\ \text{〔\%〕}}{単位なし（Ω×A÷V）}
たとえば\%$Z=5\,\%$ の変圧器は、定格の電流を流すと、中で定格の電圧の5\,\% だけ電圧が下がります。このように、量を「定格（基準）の何倍か」で表す方法を\textbf{単位法}といいます。単位法で書くと、変圧器の\qty{6600}{V}の側と\qty{210}{V}の側のように、電圧の違う所を、同じ数のままつないで計算できるのが便利な所です。

\phboxhead{ショートすると、どれだけ流れる？}
変圧器の出口のすぐ先で、何かの事故で2本の線がつながってしまった（\textbf{短絡}、ショート）とします。電流を止めるのは、変圧器の中のインピーダンス $Z$ だけです（変圧器の入口の側の電源は、とても強いとします）。流れる\textbf{短絡電流} $I_{\mathrm s}$ は、
\pheq{I_{\mathrm s}=\frac{V_{\mathrm n}}{Z}=I_{\mathrm n}\times\frac{100}{\%Z}}{A ＝ A×（単位なし）}
\label{ph:電-13:tanraku}です。\%$Z=5\,\%$ なら、定格の\textbf{20倍}の電流が流れます（\bkreidairef{9}）。\%$Z$ が小さい変圧器ほど、ふだんの電圧降下は小さくてよいのですが、ショートしたときの電流は大きくなります。

\begin{phmijika}[ブレーカーの選び方]
ショートしたときの大きな電流は、電線を熱でとかしたり、強い磁場の力で機器をこわしたりします。工場やビルでは、変圧器の\%$Z$ から短絡電流を計算し、その電流でも安全に回路を切れるブレーカー（遮断器）を選びます。
\end{phmijika}

\phkokomade

\Needspace{16\baselineskip}\phstage{shiki}{箱のくさりから、波の式を導く}{5}

\bksubsection{6}{分布定数回路：小さな箱のくさりと電信方程式}
\label{ph:電-13:sec6}
長い線では、線の抵抗やインダクタンスや静電容量が、1か所にまとまらずに、線全体にうすく広がっています。このような回路を\textbf{分布定数回路}といいます。\bkidx{ぶんぷていすうかいろ}{分布定数回路}7-11⑥では、長い線を「小さな箱のくさり」と見て、行列のかけ算で調べました。ここでは同じ考えで、線の上の電圧と電流の式を立てます。

線の長さ1 m あたりの量を、次のように決めます（記号の右上の「$'$」は「1 m あたり」の印です）。
\begin{itemize}[leftmargin=1.2\zw, itemsep=1pt, topsep=1pt, label={\color{bkTheme}●}]
\item $R'$〔Ω/m〕：行きと帰りの導体の抵抗（直列）
\item $L'$〔H/m〕：2本の導体が作るインダクタンス（直列。電-8 ⑥）
\item $G'$〔S/m〕：2本のあいだの絶縁体を通る、わずかなもれ（並列）
\item $C'$〔F/m〕：2本のあいだの静電容量（並列。電-3 ④）
\end{itemize}
長さ $\Delta x$ の短い部分は、直列の $R'\Delta x$・$L'\Delta x$ と、並列の $G'\Delta x$・$C'\Delta x$ でできた小さな箱と見なせます（図\ref{ph:電-13:fig-kusari}）。

\begin{bkfigure}
\begin{tikzpicture}[x=1mm, y=1mm]
  \draw[bkInk, line width=0.8pt] (0,16) -- (14,16) (22,16) -- (26,16) (34,16) -- (72,16) (0,0) -- (72,0);
  \sxnR{(14,16)}{8}\sxnL{(26,16)}
  \node[font=\sffamily\footnotesize, anchor=south] at (18,17.6) {$R'\Delta x$};
  \node[font=\sffamily\footnotesize, anchor=south] at (30,18.2) {$L'\Delta x$};
  \draw[bkInk, line width=0.8pt] (44,16) -- (44,12) (44,4) -- (44,0);
  \sxnRv{(44,12)}{8}
  \node[font=\sffamily\footnotesize, anchor=east] at (42.2,8) {$G'\Delta x$};
  \draw[bkInk, line width=0.8pt] (54,16) -- (54,8.7) (54,7.3) -- (54,0);
  \sxnC{(54,8)}
  \node[font=\sffamily\footnotesize, anchor=west] at (57.4,8) {$C'\Delta x$};
  \fill[bkInk] (44,16) circle (0.7) (44,0) circle (0.7) (54,16) circle (0.7) (54,0) circle (0.7);
  \draw[bkTheme, line width=0.6pt, dash pattern=on 1.2pt off 1pt, rounded corners=1mm] (8,-3) rectangle (70,24.5);
  \sxnT{(0,16)}\sxnT{(0,0)}\sxnT{(72,16)}\sxnT{(72,0)}
  \draw[bkBlue!80!black, line width=0.8pt, {Stealth[length=3pt]}-] (-3,14) -- (-3,2);
  \node[font=\sffamily\footnotesize, anchor=east] at (-4,8) {$v(x)$};
  \draw[bkBlue!80!black, line width=0.8pt, {Stealth[length=3pt]}-] (75,14) -- (75,2);
  \node[font=\sffamily\footnotesize, anchor=west] at (76,8) {$v(x+\Delta x)$};
  \draw[bkOrange!85!black, line width=1.1pt, -{Stealth[length=3pt]}] (1,19) -- (7,19);
  \node[font=\sffamily\footnotesize, anchor=south] at (4,19.4) {$i(x)$};
  \draw[bkOrange!85!black, line width=1.1pt, -{Stealth[length=3pt]}] (73,19) -- (79,19);
  \node[font=\sffamily\footnotesize, anchor=south] at (78,19.4) {$i(x+\Delta x)$};
  \node[font=\sffamily\footnotesize, anchor=north] at (39,-3.5) {長さ $\Delta x$ の箱};
  \node[anchor=west, font=\small, text width=23mm, align=left] at (96,8) {%
  この箱を、線の長さの分だけ、くさりのようにつなぐ（7-11⑥）。};
\end{tikzpicture}
\bkfigcaption{長い線の、長さ $\Delta x$ の部分を表す小さな箱。上の線が行きの導体、下の線が帰りの導体（同軸ケーブルなら外導体）。}\label{ph:電-13:fig-kusari}
\end{bkfigure}

\phboxhead{2つの式を立てる}
位置 $x$ の電圧（2本の線のあいだ）を $v(x,\ t)$、電流を $i(x,\ t)$ とします。
\begin{enumerate}[label={\bkcircled{\arabic*}}, leftmargin=2.4\zw, itemsep=2pt, topsep=2pt]
\item \textbf{電圧の下がり}（電圧則）：箱の入口と出口の電圧の差は、直列の部分にかかる電圧です（コイルの電圧は $L\,\dd i/\dd t$：電-8 ⑥）。
\[v(x)-v(x+\Delta x)=R'\Delta x\,i+L'\Delta x\,\frac{\partial i}{\partial t}\]
両辺を $\Delta x$ でわって $\Delta x\to0$ とすると（左辺は $-\partial v/\partial x$ になる：5-3）、
\[\frac{\partial v}{\partial x}=-R'i-L'\,\frac{\partial i}{\partial t}\]
\item \textbf{電流の減り}（電流則）：箱に入る電流と出る電流の差は、並列の部分に流れる電流です（コンデンサーの電流は $C\,\dd v/\dd t$：電-3 ④）。
\[i(x)-i(x+\Delta x)=G'\Delta x\,v+C'\Delta x\,\frac{\partial v}{\partial t}\]
同じようにして、
\[\frac{\partial i}{\partial x}=-G'v-C'\,\frac{\partial v}{\partial t}\]
\end{enumerate}
この2つの式の組を\textbf{電信方程式}（でんしんほうていしき）といいます。\bkidx{でんしんほうていしき}{電信方程式}遠くへ電報を送る電信のケーブルを調べるために作られたので、この名前があります（コラム）。\label{ph:電-13:denshin}

\phboxhead{数値実験：箱を100個つないだケーブル}
長さ\qty{10}{m}の\qty{50}{\ohm}のケーブル（⑦の\bkreidairef{1}の値：$L'=\qty{0.25}{\micro H/m}$、$C'=\qty{101}{pF/m}$、$R'=G'=0$）を、長さ\qty{0.1}{m}の箱100個に分け、先に\qty{75}{\ohm}の抵抗をつないで、入口に電圧の段（ステップ）を入れてみましょう。電源は、電圧 $E$ を\qty{5}{ns}かけて0から\qty{2}{V}に上げ、内部抵抗はケーブルと同じ\qty{50}{\ohm}とします。箱の電流は「電圧の差で変わる」（①）、箱の電圧は「入る電流と出る電流の差で変わる」（②）を、$\Delta t=\qty{0.01}{ns}$ ずつ進めます（6-2 のオイラー法。電流を先に新しくする形：波-1 ⑪）。はしの節は、箱の半分のコンデンサーをもつとします。
\begin{phcode}[Python]
N = 100                     \# 箱の数（1つの箱は 0.1 m）
L = 0.25e-6 * 0.1           \# 1つの箱のコイル L'Δx〔H〕
C = 101e-12 * 0.1           \# 1つの箱のコンデンサー C'Δx〔F〕
Rs, RL = 50.0, 75.0         \# 電源の内部抵抗、先の抵抗〔Ω〕
dt = 0.01e-9                \# 1歩 0.01 ns
v = [0.0] * (N + 1)         \# 節の電圧 v[0]〜v[N]
i = [0.0] * N               \# 箱のコイルの電流
for step in range(1, 12001):           \# 120 ns まで
    E = 2.0 * min(step * dt / 5e-9, 1)  \# 5 ns で 0→2 V
    for k in range(N):                 \# ①コイル：電圧の差
        i[k] += (v[k] - v[k + 1]) / L * dt
    v[0] += ((E - v[0]) / Rs - i[0]) / (C / 2) * dt
    for k in range(1, N):              \# ②コンデンサー
        v[k] += (i[k - 1] - i[k]) / C * dt
    v[N] += (i[N - 1] - v[N] / RL) / (C / 2) * dt
    if step \% 2000 == 0:               \# 20 ns ごとに表示
        print(round(step * dt * 1e9), round(v[50], 2), round(v[N], 2))
\end{phcode}
\begin{phcodeout}
20 0.0 0.0
40 1.01 0.0
60 1.0 1.19
80 1.18 1.19
100 1.2 1.2
120 1.2 1.2
\end{phcodeout}
左から、時刻〔ns〕、まん中（\qty{5}{m}）の電圧〔V〕、先の電圧〔V〕です。電源から見ると、はじめはケーブルが\qty{50}{\ohm}の抵抗のように見えるので、\qty{2}{V}は電源の内部抵抗と半分ずつに分かれ、\qty{1}{V}の段がケーブルを進みます。\qty{40}{ns}では、まん中まで来ていますが、先にはまだ着いていません。先に着くと、電圧は\qty{1.2}{V}になり、差の\qty{0.2}{V}の段がはね返って入口へ戻っていきます（図\ref{ph:電-13:fig-sim}）。\qty{1.2}{V}は、波-7 例題7 の電圧の反射係数0.20と同じ割合です（⑧）。

\begin{bkfigure}
\begin{tikzpicture}
\begin{axis}[bkplot, width=108mm, height=50mm, xmin=0, xmax=10.4, ymin=0, ymax=1.45, xtick={2,4,6,8,10}, ytick={0.5,1,1.2}, yticklabels={0.5,1.0,1.2}, xlabel={$x$（m）}, ylabel={電圧 $v$（V）}]
  \draw[bkBlue, line width=1.7pt, line join=round] (axis cs:0.0,0.998) -- (axis cs:0.2,0.997) -- (axis cs:0.4,0.997) -- (axis cs:0.6,0.997) -- (axis cs:0.8,0.998) -- (axis cs:1.0,0.998) -- (axis cs:1.2,0.998) -- (axis cs:1.4,0.998) -- (axis cs:1.6,0.997) -- (axis cs:1.8,0.998) -- (axis cs:2.0,0.998) -- (axis cs:2.2,0.996) -- (axis cs:2.4,0.997) -- (axis cs:2.6,1.002) -- (axis cs:2.8,0.993) -- (axis cs:3.0,0.997) -- (axis cs:3.2,1.004) -- (axis cs:3.4,0.987) -- (axis cs:3.6,1.01) -- (axis cs:3.8,0.985) -- (axis cs:4.0,1.007) -- (axis cs:4.2,0.997) -- (axis cs:4.4,0.981) -- (axis cs:4.6,1.017) -- (axis cs:4.8,1.042) -- (axis cs:5.0,0.92) -- (axis cs:5.2,0.774) -- (axis cs:5.4,0.588) -- (axis cs:5.6,0.339) -- (axis cs:5.8,0.14) -- (axis cs:6.0,0.042) -- (axis cs:6.2,0.009) -- (axis cs:6.4,0.002) -- (axis cs:6.6,0.0) -- (axis cs:6.8,0.0) -- (axis cs:7.0,0.0) -- (axis cs:7.2,0.0) -- (axis cs:7.4,0.0) -- (axis cs:7.6,0.0) -- (axis cs:7.8,0.0) -- (axis cs:8.0,0.0) -- (axis cs:8.2,0.0) -- (axis cs:8.4,0.0) -- (axis cs:8.6,0.0) -- (axis cs:8.8,0.0) -- (axis cs:9.0,0.0) -- (axis cs:9.2,0.0) -- (axis cs:9.4,0.0) -- (axis cs:9.6,0.0) -- (axis cs:9.8,0.0) -- (axis cs:10.0,0.0);
  \draw[bkOrange, line width=1.5pt, dash pattern=on 4pt off 2pt, line join=round] (axis cs:0.0,0.998) -- (axis cs:0.2,0.998) -- (axis cs:0.4,0.998) -- (axis cs:0.6,0.998) -- (axis cs:0.8,0.998) -- (axis cs:1.0,0.998) -- (axis cs:1.2,0.998) -- (axis cs:1.4,0.998) -- (axis cs:1.6,0.998) -- (axis cs:1.8,0.997) -- (axis cs:2.0,0.997) -- (axis cs:2.2,0.997) -- (axis cs:2.4,0.997) -- (axis cs:2.6,0.997) -- (axis cs:2.8,0.998) -- (axis cs:3.0,0.998) -- (axis cs:3.2,0.998) -- (axis cs:3.4,0.997) -- (axis cs:3.6,0.997) -- (axis cs:3.8,0.997) -- (axis cs:4.0,0.998) -- (axis cs:4.2,0.998) -- (axis cs:4.4,0.997) -- (axis cs:4.6,0.998) -- (axis cs:4.8,0.998) -- (axis cs:5.0,0.997) -- (axis cs:5.2,0.997) -- (axis cs:5.4,0.998) -- (axis cs:5.6,0.999) -- (axis cs:5.8,0.999) -- (axis cs:6.0,1.007) -- (axis cs:6.2,1.026) -- (axis cs:6.4,1.053) -- (axis cs:6.6,1.097) -- (axis cs:6.8,1.149) -- (axis cs:7.0,1.179) -- (axis cs:7.2,1.198) -- (axis cs:7.4,1.212) -- (axis cs:7.6,1.202) -- (axis cs:7.8,1.2) -- (axis cs:8.0,1.195) -- (axis cs:8.2,1.198) -- (axis cs:8.4,1.207) -- (axis cs:8.6,1.19) -- (axis cs:8.8,1.208) -- (axis cs:9.0,1.193) -- (axis cs:9.2,1.205) -- (axis cs:9.4,1.194) -- (axis cs:9.6,1.205) -- (axis cs:9.8,1.194) -- (axis cs:10.0,1.204);
  \draw[bkInk, line width=0.8pt, -{Stealth[length=2mm]}] (axis cs:5.9,0.55) -- (axis cs:6.8,0.55);
  \node[font=\sffamily\footnotesize, anchor=west, fill=white, inner sep=1pt] at (axis cs:6.9,0.55) {30 ns（実線）：右へ進む};
  \draw[bkInk, line width=0.8pt, -{Stealth[length=2mm]}] (axis cs:6.5,1.33) -- (axis cs:5.6,1.33);
  \node[font=\sffamily\footnotesize, anchor=west, fill=white, inner sep=1pt] at (axis cs:6.6,1.33) {70 ns（破線）：はね返りが左へ};
\end{axis}
\end{tikzpicture}
\bkfigcaption{数値実験の、ケーブルの上の電圧。\qty{30}{ns}（実線）では\qty{1}{V}の段が右へ進み、\qty{70}{ns}（破線）では、先ではね返った\qty{0.2}{V}の段が左へ戻ってくる。段の前の小さな波打ちは、ケーブルを箱に分けたためのもので、箱を細かくするほど小さくなる（波-5 ⑦の分散）。}\label{ph:電-13:fig-sim}
\end{bkfigure}

\noindent 箱の数を増やして箱を短くすると、結果は1つの形に近づいていきます。箱を限りなく細かくした先が、電信方程式です。電子回路を設計するシミュレーターも、長い線を、このような箱のくさりや、⑦⑨の式で計算しています。

\bksubsection{7}{損失のない線：波動方程式・速さ・特性インピーダンス}
\label{ph:電-13:sec7}
よい同軸ケーブルや、短い区間の送電線では、$R'$ と $G'$ は小さく、ほとんど0と見てかまいません。$R'=G'=0$ の線を\textbf{損失のない線}といいます。電信方程式は
\begin{bkhenkei}[0.52]
\frac{\partial v}{\partial x}=-L'\,\frac{\partial i}{\partial t} & ①で $R'=0$ \\
\frac{\partial i}{\partial x}=-C'\,\frac{\partial v}{\partial t} & ②で $G'=0$ \\
\end{bkhenkei}
となります。1つめを $x$ で微分し、2つめを入れます（微分の順番は入れかえてよい：8-13）。
\[\frac{\partial^2v}{\partial x^2}=-L'\,\frac{\partial}{\partial t}\Bigl(\frac{\partial i}{\partial x}\Bigr)=L'C'\,\frac{\partial^2v}{\partial t^2}\]
両辺を $L'C'$ でわると、
\pheq{\frac{\partial^2v}{\partial t^2}=c^2\,\frac{\partial^2v}{\partial x^2}\qquad c=\frac{1}{\sqrt{L'C'}}}{m/s（H/m×F/m＝s$^2$/m$^2$）}
\label{ph:電-13:hado}です。弦の波（波-5 ④）と同じ\textbf{波動方程式}で、線の上の電圧は、形を保って速さ $c$ で進む波になります（8-13⑨のダランベールの解）。電流 $i$ についても、同じ式が出ます。単位は、H＝V·s/A、F＝A·s/V なので H·F＝s$^2$、$1/\sqrt{\text{H/m}\times\text{F/m}}$＝m/s です。

\phboxhead{浅い水の波と同じ形}
波-5 ⑤の浅い水の波では、水面の高さ $\eta$ と水の速さ $U$ の2つの式から波動方程式を作りました。電信方程式（損失なし）と並べると、形がそっくりです。
\begin{bkcenter}
\begin{bktable}{colspec={Q[c,m,wd=37mm] Q[c,m,wd=37mm] Q[l,m,wd=42mm]}, rowsep=1.5pt, row{1}={bg=bkTheme!10, font=\sffamily\small}}
浅い水の波（波-5 ⑤） & 電線の波（この章） & 対応 \\
$\dfrac{\partial\eta}{\partial t}=-h\,\dfrac{\partial U}{\partial x}$ & $\dfrac{\partial v}{\partial t}=-\dfrac{1}{C'}\,\dfrac{\partial i}{\partial x}$ & 高さ $\eta$ と 電圧 $v$ \\
$\dfrac{\partial U}{\partial t}=-g\,\dfrac{\partial\eta}{\partial x}$ & $\dfrac{\partial i}{\partial t}=-\dfrac{1}{L'}\,\dfrac{\partial v}{\partial x}$ & 速さ $U$ と 電流 $i$ \\
$c=\sqrt{gh}$ & $c=\sqrt{\dfrac{1}{L'}\times\dfrac{1}{C'}}$ & $h$ と $1/C'$、$g$ と $1/L'$ \\
\end{bktable}
\bktabcaption{浅い水の波と、損失のない電線の波。どちらも「たまる量の変わり方は、流れの場所による違いで決まる」（1行め：水の量・電荷の保存）と「流れの変わり方は、押すものの場所による違いで決まる」（2行め：運動方程式・電圧則）の組。}\label{ph:電-13:tab-mizu}
\end{bkcenter}

\phboxhead{特性インピーダンス}
右へ進む電圧の波 $v=v_0(x-ct)$（$v_0$ は形を表す関数）のとき、電流はどうなるでしょう。電流も同じ形で $i=v_0(x-ct)/Z$ と書けると考えて、1つめの式 $\partial v/\partial x=-L'\,\partial i/\partial t$ に入れます。$v_0$ の微分を $v_0'$ と書くと、$\partial v/\partial x=v_0'$、$\partial i/\partial t=-c\,v_0'/Z$（合成関数の微分：5-6）なので、
\[v_0'=L'c\,\frac{v_0'}{Z}\qquad\text{だから}\qquad Z=L'c=\frac{L'}{\sqrt{L'C'}}=\sqrt{\frac{L'}{C'}}\]
です。右へ進む波では、電圧と電流の比がいつも同じ値になります。これが、①で見た\textbf{特性インピーダンス} $Z_0$ です。
\pheq{Z_0=\frac{v}{i}=\sqrt{\frac{L'}{C'}}}{Ω ＝ $\sqrt{\text{H/F}}$}
\label{ph:電-13:Z0}H/F＝(V·s/A)÷(A·s/V)＝V$^2$/A$^2$ なので、$\sqrt{\text{H/F}}$＝Ω です。左へ進む波 $v=v_0(x+ct)$ では、同じ計算で $i=-v/Z_0$ になります。電流の向き（右向きを正）が逆になるからです。これは、波-7 ③で、左へ進む波の力が $-Zu$ だったのと同じです。

\begin{phyomi}[$Z_0=\sqrt{L'/C'}$、$c=1/\sqrt{L'C'}$]
「$L'$（電流の変わりにくさ）が大きいほど、同じ電圧で流れる電流は小さく（$Z_0$ が大きい）、$C'$（電荷のためやすさ）が大きいほど、電流は大きい（$Z_0$ が小さい）。速さは、$L'$ と $C'$ のどちらが大きくてもおそくなる」。弦の $Z=\sqrt{F_{\mathrm T}\mu}$・$c=\sqrt{F_{\mathrm T}/\mu}$（波-7 ③）と同じ形です。$Z_0$ は線の長さによらない、線の形と材料だけで決まる値です。テスターではかる線の抵抗（電-4）とは、まったく別の量です。
\end{phyomi}

\begin{phflow}[損失のない線]
\phflowrow{$\dfrac{\partial v}{\partial x}=-R'i-L'\dfrac{\partial i}{\partial t}$\\$\dfrac{\partial i}{\partial x}=-G'v-C'\dfrac{\partial v}{\partial t}$}{$R'=G'=0$\\（損失のない線）}{$\dfrac{\partial^2v}{\partial t^2}=\dfrac{1}{L'C'}\dfrac{\partial^2v}{\partial x^2}$}
\phflowrow{$\dfrac{\partial^2v}{\partial t^2}=\dfrac{1}{L'C'}\dfrac{\partial^2v}{\partial x^2}$}{右へ進む波\\$v=v_0(x-ct)$}{$c=\dfrac{1}{\sqrt{L'C'}}$、$Z_0=\sqrt{\dfrac{L'}{C'}}$}
\end{phflow}

\phboxhead{同軸ケーブルの $L'$・$C'$}
内導体の半径を $a$、外導体の内側の半径を $b$、あいだを比誘電率 $\varepsilon_{\mathrm r}$ の絶縁体で満たした同軸ケーブルでは、
\pheq{C'=\frac{2\pi\varepsilon_{\mathrm r}\varepsilon_0}{\ln(b/a)}\qquad L'=\frac{\mu_0}{2\pi}\,\ln\frac{b}{a}}{F/m・H/m}
です。$C'$ は電-3 ⑩（ガウスの法則で、あいだの電場を $\lambda/(2\pi\varepsilon r)$ と求め、内から外まで積分して電圧を出した：$\lambda$ は1 m あたりの電荷）、$L'$ は電-8 ⑥（アンペールの法則で、あいだの磁場を $\mu_0I/(2\pi r)$ と求め、1 m あたりの磁束を数えた：電-6 ⑤）で導きました。どちらにも $\ln(b/a)$ が出るのは、場が $1/r$ に比例し、それを $r$ で積分するからです。これを⑦の式に入れると、
\begin{bkhenkei}[0.52]
c=\frac{1}{\sqrt{L'C'}}=\frac{1}{\sqrt{\mu_0\varepsilon_{\mathrm r}\varepsilon_0}}=\frac{c_0}{\sqrt{\varepsilon_{\mathrm r}}} & $\ln(b/a)$ と $2\pi$ が消える \\
Z_0=\sqrt{\frac{L'}{C'}}=\frac{1}{2\pi}\sqrt{\frac{\mu_0}{\varepsilon_{\mathrm r}\varepsilon_0}}\,\ln\frac{b}{a} & $\ln(b/a)$ の1つは分子に残る \\
\approx\frac{60}{\sqrt{\varepsilon_{\mathrm r}}}\,\ln\frac{b}{a}\quad\text{〔Ω〕} & $\dfrac{1}{2\pi}\sqrt{\dfrac{\mu_0}{\varepsilon_0}}\approx\qty{60}{\ohm}$ \\
\end{bkhenkei}
\label{ph:電-13:coax}です。ここで $c_0=1/\sqrt{\mu_0\varepsilon_0}$ は真空の光の速さです（電-15 ⑥で導く）。\textbf{ケーブルの中の速さは、太さによらず、絶縁体の比誘電率だけで決まり}、ポリエチレン（$\varepsilon_{\mathrm r}\approx2.3$）では $c_0$ の約0.66倍です（\bkreidairef{1}）。特性インピーダンスは、太さの\textbf{比} $b/a$ で決まります。ケーブル全体を2倍に太くしても、$Z_0$ は変わりません。

\begin{phbaisuru}{太さの比 $b/a$}
$\ln(b/a)$ は $\ln2\approx0.69$ だけ増えるので、$Z_0$ は2倍にはならず、ポリエチレンなら約\qty{27}{\ohm}増えます。$Z_0$ は比の対数で決まるので、\qty{50}{\ohm}や\qty{75}{\ohm}にするには、比を3〜5くらいにすればよく、ケーブルの形として作りやすい値です。
\par\smallskip\noindent\begin{tikzpicture}
\begin{axis}[bkplot, width=100mm, height=48mm, xmin=1, xmax=8.4, ymin=0, ymax=105, xtick={1,2,3,4,5,6,7,8}, ytick={25,50,75,100}, xlabel={太さの比 $b/a$}, ylabel={$Z_0$（Ω）}]
  \addplot[bkcurve, domain=1:8.2, samples=60] {60/sqrt(2.3)*ln(x)};
  \addplot[bkOrange, line width=1.5pt, dash pattern=on 4pt off 2pt, domain=1:8.2, samples=60] {60/sqrt(1.5)*ln(x)};
  \fill[bkRed] (axis cs:3.56,50.2) circle (1.8pt) (axis cs:4.62,75) circle (1.8pt);
  \node[font=\sffamily\footnotesize, anchor=north west, inner sep=1pt] at (axis cs:3.65,48) {50 Ω};
  \node[font=\sffamily\footnotesize, anchor=south east, inner sep=1pt] at (axis cs:4.5,77) {75 Ω};
\end{axis}
\end{tikzpicture}
\par\noindent（実線：ポリエチレン $\varepsilon_{\mathrm r}=2.3$、破線：発泡ポリエチレン $\varepsilon_{\mathrm r}=1.5$。●は例題1と例題2の値。）
\end{phbaisuru}

\begin{phbaisuru}{1 m あたりのインダクタンス $L'$}
$Z_0=\sqrt{L'/C'}$ は $\sqrt2\approx1.41$ 倍、速さ $c=1/\sqrt{L'C'}$ は $1/\sqrt2\approx0.71$ 倍になります。$C'$ を2倍にすると、$Z_0$ も速さも $1/\sqrt2$ 倍です。線の長さを2倍にすると、$Z_0$ と速さは変わらず、パルスが往復する時間が2倍になります。
\end{phbaisuru}

\begin{phmodel}[同軸ケーブルの $L'$・$C'$]
\phignore{導線の抵抗（$R'$）と絶縁体のもれ（$G'$）。内導体の中の磁場（周波数が高いと、電流は表面だけを流れる：電-8 ⑫の表皮効果）。}
\phassume{内導体と外導体はまん丸で、中心が同じ。絶縁体は磁石につかない（透磁率 $\mu_0$）。}
\phlimit{周波数が低く、電流が導線の中まで流れるとき（$L'$ が少し大きくなる）。とても高い周波数で、$R'$ による弱まりが大きいとき（⑨）。}
\end{phmodel}

\bksubsection{8}{はね返りと整合}
\label{ph:電-13:sec8}
特性インピーダンス $Z_0$ の線の先（$x=l$）に、インピーダンス $Z_{\text{先}}$ の機器（アンテナ、テレビ、抵抗）をつなぎます。右から来た波（電圧 $v_{\text{来}}$、電流 $v_{\text{来}}/Z_0$）が先に着くと、一部がはね返ります。はね返った波の電圧を $v_{\text{返}}$ とすると、左へ進む波なので、その電流は $-v_{\text{返}}/Z_0$ です（⑦）。

\phboxhead{先で守る2つの条件}
先の所では、線の電圧と電流が、そのまま機器の電圧と電流です。
\begin{enumerate}[label={\bkcircled{\arabic*}}, leftmargin=2.4\zw, itemsep=2pt, topsep=2pt]
\item 電圧は、来た波とはね返った波の和：$v=v_{\text{来}}+v_{\text{返}}$
\item 電流も和：$i=\dfrac{v_{\text{来}}}{Z_0}-\dfrac{v_{\text{返}}}{Z_0}$。そして機器では $v=Z_{\text{先}}\,i$
\end{enumerate}
①を②に入れると、$v_{\text{来}}+v_{\text{返}}=\dfrac{Z_{\text{先}}}{Z_0}(v_{\text{来}}-v_{\text{返}})$。両辺に $Z_0$ をかけて、$v_{\text{返}}$ を左に集めると、$(Z_0+Z_{\text{先}})v_{\text{返}}=(Z_{\text{先}}-Z_0)v_{\text{来}}$ なので、
\pheq{r_V=\frac{v_{\text{返}}}{v_{\text{来}}}=\frac{Z_{\text{先}}-Z_0}{Z_{\text{先}}+Z_0}}{単位なし}
\label{ph:電-13:rV}です。$r_V$ を\textbf{電圧の反射係数}といいます。\bkidx{はんしゃけいすう（でんあつ）}{反射係数（電圧）}

\phboxhead{波-7 の $r_{\text{押}}$ と同じ式}
波-7 ④では、反射係数に2つの数え方がありました。速さ（流れるもの）で数えた $r$ と、力・音圧（押すもの）で数えた $r_{\text{押}}=-r$ です（思い出そう）。電気の線では、\textbf{電圧が押すもの、電流が流れるもの}です。線を「物1」（$Z_1=Z_0$）、先の機器を「物2」（$Z_2=Z_{\text{先}}$）と見ると、
\begin{itemize}[leftmargin=1.2\zw, itemsep=1pt, topsep=1pt, label={\color{bkTheme}●}]
\item 電圧の反射係数 $r_V=\dfrac{Z_{\text{先}}-Z_0}{Z_{\text{先}}+Z_0}$ は、波-7 の $r_{\text{押}}$ と同じ式（押すもので数えた）。
\item 電流の反射係数 $r_I=\dfrac{-v_{\text{返}}/Z_0}{v_{\text{来}}/Z_0}=-r_V=\dfrac{Z_0-Z_{\text{先}}}{Z_0+Z_{\text{先}}}$ は、波-7 の $r$ と同じ式（流れるもので数えた）。
\end{itemize}
電気の本は、ふつう電圧の反射係数 $r_V$ を「反射係数」とよびます（音響の本が音圧の $r_{\text{押}}$ を使うのと同じ）。はね返る電力の割合は、どちらで数えても $r_V^{\,2}=r_I^{\,2}$ です。先の機器へ伝わる電圧は $v=(1+r_V)v_{\text{来}}$ で、伝わる電力の割合は
\[\frac{(1+r_V)^2v_{\text{来}}^{\,2}/Z_{\text{先}}}{v_{\text{来}}^{\,2}/Z_0}=\frac{4Z_0Z_{\text{先}}}{(Z_0+Z_{\text{先}})^2}=1-r_V^{\,2}\]
です（波-7 ④の「通りぬける割合」と同じ式。右へ進む波の電力は $vi=v^2/Z_0$）。

\begin{phflow}[電圧の反射係数]
\phflowrow{$r_V=\dfrac{Z_{\text{先}}-Z_0}{Z_{\text{先}}+Z_0}$}{$Z_{\text{先}}=Z_0$\\（そろえる）}{$r_V=0$（整合）}
\phflowrow{$r_V=\dfrac{Z_{\text{先}}-Z_0}{Z_{\text{先}}+Z_0}$}{$Z_{\text{先}}\to\infty$\\（先を開ける）}{$r_V=+1$}
\phflowrow{$r_V=\dfrac{Z_{\text{先}}-Z_0}{Z_{\text{先}}+Z_0}$}{$Z_{\text{先}}=0$\\（先をショート）}{$r_V=-1$}
\end{phflow}
\label{ph:電-13:seigou}①の図\ref{ph:電-13:fig-tdr}の3つの画面は、この3つの場合でした。先を開けると、電流は0（$r_I=-1$：なわの固定端と同じ形）で、電圧は2倍に重なります。ショートさせると、電圧は0（$r_V=-1$）で、電流が2倍に重なります。$Z_{\text{先}}=Z_0$ にそろえて、はね返りをなくすことを\textbf{整合}（マッチング）といいます（波-7 ⑧）。\bkidx{せいごう}{整合}

\phboxhead{時間の図で見る}
⑥の数値実験（\qty{50}{\ohm}のケーブルの先に\qty{75}{\ohm}）を、横に位置、下へ時刻をとった図にかくと、図\ref{ph:電-13:fig-jikan}のようになります。$r_V=(75-50)/(75+50)=0.20$ なので、\qty{1}{V}の段が先に着くと、\qty{0.2}{V}の段がはね返り、先の電圧は $1+0.2=\qty{1.2}{V}$ になります。

\begin{bkfigure}
\begin{tikzpicture}[x=1mm, y=1mm]
  \fill[bkBlue!8] (0,0) -- (60,-30) -- (60,-62) -- (0,-62) -- cycle;
  \fill[bkOrange!15] (0,-60) -- (60,-30) -- (60,-62) -- (0,-62) -- cycle;
  \draw[bkInk, line width=0.6pt, -{Stealth[length=2mm]}] (0,0) -- (66,0);
  \draw[bkInk, line width=0.6pt, -{Stealth[length=2mm]}] (0,0) -- (0,-66);
  \draw[bkInk, line width=0.6pt] (60,1) -- (60,-62);
  \node[font=\sffamily\footnotesize, anchor=south] at (0,1) {入口};
  \node[font=\sffamily\footnotesize, anchor=south] at (60,1) {先（75 Ω）};
  \node[font=\sffamily\footnotesize, anchor=west] at (67,0) {$x$};
  \node[font=\sffamily\footnotesize, anchor=north] at (0,-66.5) {時刻 $t$};
  \draw[bkBlue!80!black, line width=1.3pt, -{Stealth[length=2.4mm]}] (0,0) -- (60,-30);
  \draw[bkRed!80!black, line width=1.3pt, -{Stealth[length=2.4mm]}] (60,-30) -- (0,-60);
  \node[font=\sffamily\footnotesize, anchor=east] at (-1,-30) {50 ns};
  \node[font=\sffamily\footnotesize, anchor=east] at (-1,-60) {101 ns};
  \draw[bkGray, line width=0.4pt, dashed] (0,-30) -- (60,-30);
  \node[font=\sffamily\footnotesize, anchor=south west, fill=white, inner sep=1pt] at (30,-14) {来た段 1 V};
  \node[font=\sffamily\footnotesize, anchor=north west, inner sep=1pt, align=left] at (31,-46.5) {はね返った\\段 0.2 V};
  \node[font=\sffamily\footnotesize, fill=white, inner sep=1pt] at (40,-5) {まだ 0 V};
  \node[font=\sffamily\footnotesize, fill=bkBlue!8, inner sep=1pt] at (17,-24) {1.0 V};
  \node[font=\sffamily\footnotesize, inner sep=1pt] at (52,-58.5) {1.2 V};
  \node[anchor=west, font=\small, text width=50mm, align=left] at (72,-30) {%
  線がまっすぐなのは、段が一定の速さ $c$ で進むから。傾きは $1/c$（1 m 進むのに約\qty{5}{ns}）。\par\vspace{3pt}
  入口の電源の内部抵抗は\qty{50}{\ohm}で、線とそろっているので、戻ってきた段は入口でははね返らない。};
\end{tikzpicture}
\bkfigcaption{\qty{10}{m}・\qty{50}{\ohm}のケーブルの先に\qty{75}{\ohm}をつないだときの、時間の図（下へ行くほど後の時刻）。色の所は、その場所・時刻の電圧。}\label{ph:電-13:fig-jikan}
\end{bkfigure}

\phboxhead{正弦波では、定在波ができる}
正弦波の信号を送り続けると、来た波とはね返った波が重なって、線の上に\textbf{定在波}（波-6 ③）ができます。電圧の振幅は、2つの波の山が重なる所で $(1+|r_V|)$ 倍、山と谷が重なる所で $(1-|r_V|)$ 倍になり、その間を半波長ごとに繰り返します（図\ref{ph:電-13:fig-teizai}）。いちばん大きい所といちばん小さい所の比
\pheq{\text{SWR}=\frac{1+|r_V|}{1-|r_V|}}{単位なし}
を\textbf{定在波比}（電圧定在波比、SWR）といいます。\bkidx{ていざいはひ}{定在波比}整合していれば $\text{SWR}=1$ で、線の上のどこでも振幅は同じです。\qty{50}{\ohm}の線に\qty{75}{\ohm}をつなぐと、$(1+0.2)/(1-0.2)=1.5$ です。アマチュア無線の人は、SWR をはかる計器で、アンテナとケーブルがそろっているかを確かめます。

\begin{bkfigure}
\begin{tikzpicture}
\begin{axis}[bkplot, width=100mm, height=48mm, xmin=0, xmax=1.02, ymin=0, ymax=2.3, xtick={0,0.25,0.5,0.75,1}, xticklabels={0,$\lambda/4$,$\lambda/2$,$3\lambda/4$,$\lambda$}, ytick={0.8,1.2,2}, yticklabels={0.8,1.2,2.0}, xlabel={先からの距離}, ylabel={電圧の振幅（来た波の何倍）}]
  \addplot[bkcurve, domain=0:1, samples=200] {sqrt(1+0.04+0.4*cos(720*x))};
  \addplot[bkOrange, line width=1.5pt, dash pattern=on 4pt off 2pt, domain=0:1, samples=200] {2*abs(cos(360*x))};
\end{axis}
\end{tikzpicture}
\bkfigcaption{線の上の電圧の振幅（実線：$r_V=0.2$、破線：先を開けたとき）。横軸は線の先からもどった距離で、左はし（0）が線の先。振幅は半波長ごとに同じ形を繰り返す。先を開けると、先が腹、$\lambda/4$ もどった所が節になる（波-6 の閉じた管と同じ形）。}\label{ph:電-13:fig-teizai}
\end{bkfigure}

\phboxhead{はてな？の答え}
\label{ph:電-13:hatena}\bkhatenaref の答えです。テレビの電波（\qtyrange{470}{710}{MHz}）の信号は、ケーブルの中では波長が数十 cm しかありません。屋根からテレビまでの\qty{10}{m}ほどのケーブルは、波長の何十倍もの「長い線」です。電気の信号は、線の上を波として進み、特性インピーダンスの違う所があると、そこではね返ります。

ふつうの2本の電線を、間かくを決めずに並べると、場所によって間かくや、周りの物との距離が変わり、特性インピーダンスがあちこちで変わります。そのたびに波の一部がはね返り、はね返った波は少しおくれてテレビに届きます。昔のアナログ放送では、これが画面の像が二重に見える「\textbf{ゴースト}」の原因の1つでした（今のデジタル放送では、画面がモザイクのように乱れたり、映らなくなったりします）。また、電場と磁場が外の空間に広がっているので、周りの雑音の電波を拾いやすく、信号も外へもれます。

同軸ケーブルは、太さの比 $b/a$ と絶縁体で、特性インピーダンスがケーブル全体でぴったり\qty{75}{\ohm}になるように作ってあります（\bkreidairef{2}）。テレビの入口も\qty{75}{\ohm}にそろえてあるので（整合）、はね返りがほとんどありません。そのうえ、電場と磁場は外導体の内側にとじこめられ、外の雑音の電波もさえぎられます。\textbf{特性インピーダンスが決まっていて、はね返りと雑音を防げる}から、テレビのアンテナ線には同軸ケーブルを使うのです。

\bksubsection{9}{損失のある線：双曲線関数の形}
\label{ph:電-13:sec9}
本物の線では、$R'$ と $G'$ が0ではなく、波は進むうちに弱まります。周波数 $f$（角周波数 $\omega$）の正弦波だけを考えると、電圧と電流を複素数（フェーザ）$V(x)$・$I(x)$ で書けます（電-10 ③）。時刻の微分は $j\omega$ をかけることなので（6-10⑩）、電信方程式は
\begin{bkhenkei}[0.5]
\frac{\dd V}{\dd x}=-(R'+j\omega L')I=-Z'I & ①。$Z'=R'+j\omega L'$（1 m あたりの直列のインピーダンス） \\
\frac{\dd I}{\dd x}=-(G'+j\omega C')V=-Y'V & ②。$Y'=G'+j\omega C'$（1 m あたりの並列のアドミタンス） \\
\end{bkhenkei}
になります。7-11⑥の $z$・$y$ が、この $Z'$・$Y'$ です。①を $x$ で微分して②を入れると、
\[\frac{\dd^2V}{\dd x^2}=-Z'\,\frac{\dd I}{\dd x}=Z'Y'\,V=\gamma^2V\qquad\gamma=\sqrt{Z'Y'}\]
で、6-11⑤の $y''=a^2y$ と同じ形です。$\gamma$ を\textbf{伝搬定数}（でんぱんていすう）といいます。\bkidx{でんぱんていすう}{伝搬定数}答えは $\cosh\gamma x$ と $\sinh\gamma x$ の組み合わせ（$e^{-\gamma x}$ と $e^{\gamma x}$ の組み合わせと同じ）です。

\phboxhead{長さ $l$ の線の四端子定数}
線の先（$x=l$）の電圧・電流を $V_2$・$I_2$ として、次の形を考えます。
\[V(x)=V_2\cosh\gamma(l-x)+Z_0I_2\sinh\gamma(l-x)\]
\[I(x)=\frac{V_2}{Z_0}\sinh\gamma(l-x)+I_2\cosh\gamma(l-x)\]
ここで $Z_0=\sqrt{Z'/Y'}$ です。$x=l$ では $\cosh0=1$、$\sinh0=0$ なので、ちゃんと $V_2$・$I_2$ になります。$x$ で微分すると（$\cosh$ の微分は $\sinh$、$\sinh$ の微分は $\cosh$：6-11④。中身の $l-x$ の微分 $-1$ がつく）、
\[\frac{\dd V}{\dd x}=-\gamma\bigl(V_2\sinh\gamma(l-x)+Z_0I_2\cosh\gamma(l-x)\bigr)=-\gamma Z_0\,I(x)\]
で、$\gamma Z_0=\sqrt{Z'Y'}\sqrt{Z'/Y'}=Z'$ なので①が成り立ちます。②も同じように確かめられます（章末問題）。入口 $x=0$ の値を行列で書くと、
\pheq{\begin{pmatrix}V_1\\I_1\end{pmatrix}=\begin{pmatrix}\cosh\gamma l&Z_0\sinh\gamma l\\[2pt]\dfrac{1}{Z_0}\sinh\gamma l&\cosh\gamma l\end{pmatrix}\begin{pmatrix}V_2\\I_2\end{pmatrix}}{$\gamma l$ は単位なし}
\label{ph:電-13:cosh}で、7-11⑥・9-11⑨と同じ行列です。$AD-BC=\cosh^2\gamma l-\sinh^2\gamma l=1$（6-11③）で、③の約束どおりになっています。

\phboxhead{伝搬定数を2つに分ける}
$\gamma$ は複素数なので、実部と虚部に分けて $\gamma=\alpha+j\beta$ と書きます。右へ進む波は $e^{-\gamma x}=e^{-\alpha x}e^{-j\beta x}$ の形です。
\begin{itemize}[leftmargin=1.2\zw, itemsep=1pt, topsep=1pt, label={\color{bkTheme}●}]
\item $\alpha$〔1/m〕（\textbf{減衰定数}）：1 m 進むごとに、振幅が $e^{-\alpha}$ 倍に弱まる（波-7 ⑦の吸収の定数と同じ）。
\item $\beta$〔rad/m〕（\textbf{位相定数}）：1 m 進むごとに、位相が $\beta$ だけずれる。波-5 ③の波数 $k$ にあたり、波長は $\lambda=2\pi/\beta$。
\end{itemize}
損失のない線（$R'=G'=0$）では、$\gamma=\sqrt{j\omega L'\times j\omega C'}=j\omega\sqrt{L'C'}$ で、$\alpha=0$、$\beta=\omega\sqrt{L'C'}=\omega/c$ です。$Z_0=\sqrt{j\omega L'/(j\omega C')}=\sqrt{L'/C'}$ で、⑦と同じになります。

\begin{phflow}[長さ $l$ の線の行列]
\phflowrow{$\begin{pmatrix}\cosh\gamma l&Z_0\sinh\gamma l\\\frac{1}{Z_0}\sinh\gamma l&\cosh\gamma l\end{pmatrix}$}{損失なし\\$\gamma=j\beta$}{$\begin{pmatrix}\cos\beta l&jZ_0\sin\beta l\\\frac{j}{Z_0}\sin\beta l&\cos\beta l\end{pmatrix}$}
\phflowrow{$\begin{pmatrix}\cosh\gamma l&Z_0\sinh\gamma l\\\frac{1}{Z_0}\sinh\gamma l&\cosh\gamma l\end{pmatrix}$}{直流（$\omega=0$）\\$\gamma=\sqrt{R'G'}$}{7-11⑥の\\$\cosh1$・$\sinh1$}
\end{phflow}
1段めは、$\cosh(j\theta)=\cos\theta$、$\sinh(j\theta)=j\sin\theta$（6-11 と 6-10⑥の $e^{j\theta}$ から）を使いました。損失のない線の行列は、⑩・⑪で使います。

\begin{motto}[5]
\textbf{ひずまない線}：$\alpha$ と $\beta$ は、ふつう周波数によって変わります。周波数で弱まり方や速さが違うと、いろいろな周波数をふくむパルスは、進むうちに形がくずれます（ひずみ）。ところが、$R'/L'=G'/C'$ が成り立つ線では、計算すると $\alpha=\sqrt{R'G'}$、$\beta=\omega\sqrt{L'C'}$ になり、弱まり方はどの周波数でも同じ、速さも同じで、形がくずれずに、大きさだけ小さくなって届きます。これを\textbf{ひずみのない条件}（ヘビサイドの条件）といいます。ふつうの電話線では $R'/L'$ のほうがずっと大きいので、線のとちゅうにコイルを入れて $L'$ を大きくする方法（装荷）が、1900年ごろから長い電話線に使われました。
\end{motto}

\bksubsection{10}{長い送電線とフェランチ効果}
\label{ph:電-13:sec10}
山の向こうの発電所から都市まで、何百 km もの送電線で電気を送るとき、\qty{50}{Hz}でも、線は「短い」とは言えなくなります。空気の中の架空送電線では、線の上の速さはほぼ $c_0$ で、波長は\qty{6000}{km}です。\qty{300}{km}の線は、波長の20分の1です。

$R'$・$G'$ を無視して、損失のない線の行列（⑨）を使います。先に何もつながない（負荷を切った）とき、$I_2=0$ なので、
\[V_1=\cos\beta l\times V_2\qquad\text{だから}\qquad V_2=\frac{V_1}{\cos\beta l}\]
です。$\cos\beta l<1$ なので、\textbf{受電端の電圧 $V_2$ は、送電端の電圧 $V_1$ より高くなります}。これを\textbf{フェランチ効果}といいます。\bkidx{ふぇらんちこうか}{フェランチ効果}\label{ph:電-13:ferranti}\qty{300}{km}では $\beta l=2\pi\times\tfrac{300}{6000}=\tfrac{\pi}{10}$ で、$1/\cos(\pi/10)\approx1.05$、つまり5\,\% 高くなります（\bkreidairef{10}）。

\begin{bkfigure}
\begin{tikzpicture}
\begin{axis}[bkplot, width=104mm, height=46mm, xmin=0, xmax=1050, ymin=0.9, ymax=2.15, xtick={200,400,600,800,1000}, ytick={1,1.24,2}, yticklabels={1.00,1.24,2.00}, xlabel={線の長さ $l$（km）}, ylabel={$V_2/V_1$}]
  \addplot[bkcurve, domain=0:1000, samples=100] {1/cos(0.06*x)};
  \fill[bkRed] (axis cs:300,1.0515) circle (1.8pt) (axis cs:600,1.236) circle (1.8pt) (axis cs:1000,2) circle (1.8pt);
  \draw[bkGray, line width=0.4pt, dashed] (axis cs:300,0.9) -- (axis cs:300,1.0515) (axis cs:600,0.9) -- (axis cs:600,1.236);
  \node[font=\sffamily\footnotesize, anchor=west, inner sep=1pt] at (axis cs:40,1.85) {300 km で1.05倍、600 km で1.24倍、1000 km で2倍};
\end{axis}
\end{tikzpicture}
\bkfigcaption{\qty{50}{Hz}・損失のない線で、先を開けたときの受電端と送電端の電圧の比 $1/\cos\beta l$。長さ\qty{1500}{km}（4分の1波長）で限りなく大きくなる。}\label{ph:電-13:fig-ferranti}
\end{bkfigure}

\begin{phyomi}[$V_2=V_1/\cos\beta l$]
「線の静電容量を充電する電流が、線のインダクタンスを通ると、先の電圧を押し上げる」。負荷がなくても、線の $C'$ には交流の充電電流が流れます。コンデンサーの電流は電圧より90°進み、それがコイルを通ると、電圧が同じ向きに加わるので（④の図のリアクタンスの矢印が、こんどは $V_{\mathrm r}$ と同じ向きになる）、先へ行くほど電圧が上がるのです。線の長さが4分の1波長に近づくと、線全体が共振（電-10 ⑦）に近づきます。
\end{phyomi}

\phboxhead{3つのモデル：短い線・中くらいの線・長い線}
電力の工学では、送電線の長さによって、3つのモデルを使い分けます（図\ref{ph:電-13:fig-model3}）。線全体の直列のインピーダンスを $Z=(R'+j\omega L')l$、並列のアドミタンスを $Y=(G'+j\omega C')l$ とします。

\begin{bkfigure}
\begin{tikzpicture}[x=1mm, y=1mm]
  % 短い線
  \draw[bkInk, line width=0.8pt] (0,12) -- (8,12) (16,12) -- (24,12) (0,0) -- (24,0);
  \sxnR{(8,12)}{8}
  \node[font=\sffamily\footnotesize, anchor=south] at (12,13.6) {$Z$};
  \sxnT{(0,12)}\sxnT{(0,0)}\sxnT{(24,12)}\sxnT{(24,0)}
  \node[font=\sffamily\footnotesize, anchor=north, align=center] at (12,-2) {短い線\\（④）};
  % π形
  \begin{scope}[shift={(36,0)}]
  \draw[bkInk, line width=0.8pt] (0,12) -- (14,12) (22,12) -- (36,12) (0,0) -- (36,0);
  \sxnR{(14,12)}{8}
  \node[font=\sffamily\footnotesize, anchor=south] at (18,13.6) {$Z$};
  \draw[bkInk, line width=0.8pt] (6,12) -- (6,9.5) (6,2.5) -- (6,0) (30,12) -- (30,9.5) (30,2.5) -- (30,0);
  \sxnRv{(6,9.5)}{7}\sxnRv{(30,9.5)}{7}
  \node[font=\sffamily\footnotesize, anchor=east] at (4.4,6) {$Y/2$};
  \node[font=\sffamily\footnotesize, anchor=west] at (31.8,6) {$Y/2$};
  \fill[bkInk] (6,12) circle (0.6) (6,0) circle (0.6) (30,12) circle (0.6) (30,0) circle (0.6);
  \sxnT{(0,12)}\sxnT{(0,0)}\sxnT{(36,12)}\sxnT{(36,0)}
  \node[font=\sffamily\footnotesize, anchor=north, align=center] at (18,-2) {中くらいの線\\（π形）};
  \end{scope}
  % 長い線
  \begin{scope}[shift={(86,0)}]
  \draw[bkInk, line width=2.0pt] (0,12) -- (36,12);
  \draw[bkGray!70, line width=1.0pt] (0,12) -- (36,12);
  \draw[bkInk, line width=0.8pt] (0,0) -- (36,0);
  \foreach \xx in {4,10,...,34} {\draw[bkInk!50, line width=0.4pt] (\xx,10.6) -- (\xx,1.4);}
  \sxnT{(0,12)}\sxnT{(0,0)}\sxnT{(36,12)}\sxnT{(36,0)}
  \node[font=\sffamily\footnotesize, anchor=north, align=center] at (18,-2) {長い線\\（⑨の $\cosh$・$\sinh$）};
  \end{scope}
\end{tikzpicture}
\bkfigcaption{送電線の3つのモデル。短い線は直列の $Z$ だけ、中くらいの線は並列の $Y$ を半分ずつ両はしに置いた π（パイ）形、長い線は細かく分けた箱のくさり（⑥）として、⑨の行列で計算する。}\label{ph:電-13:fig-model3}
\end{bkfigure}

\noindent π形の行列は、並列 $Y/2$・直列 $Z$・並列 $Y/2$ の3つの箱の積（③）で、
\[\begin{pmatrix}1&0\\\frac{Y}{2}&1\end{pmatrix}\begin{pmatrix}1&Z\\0&1\end{pmatrix}\begin{pmatrix}1&0\\\frac{Y}{2}&1\end{pmatrix}=\begin{pmatrix}1+\frac{ZY}{2}&Z\\Y\bigl(1+\frac{ZY}{4}\bigr)&1+\frac{ZY}{2}\end{pmatrix}\]
です。損失のない\qty{300}{km}の線では、$ZY=j\omega L'l\times j\omega C'l=-(\beta l)^2$ なので、先を開けたときの電圧の比は
\[\frac{V_2}{V_1}=\frac{1}{1+ZY/2}=\frac{1}{1-(\beta l)^2/2}=\frac{1}{1-0.0493}\approx1.052\]
で、正しい値 $1/\cos\beta l\approx1.051$ にとても近くなります（$\cos\theta\approx1-\theta^2/2$：5-17）。短い線のモデルでは並列の $Y$ がないので、電圧は上がらず、1.00 のままです。\qty{600}{km}では、π形が1.246、正しい値が1.236 と差が開いてきます。線が長くなるほど、⑨の $\cosh$・$\sinh$ の行列が必要になるのです。

\begin{phmijika}[夜の電圧の上がりすぎを防ぐ]
長い送電線や、静電容量の大きい地中ケーブル（同軸に近い形で、$C'$ が架空線の何十倍もある）では、夜など使う電気が少ないときに、受電端の電圧が上がりすぎることがあります。変電所では、コイル（分路リアクトル）を線と大地のあいだにつないで、線の静電容量の電流を打ち消し、電圧を抑えています。
\end{phmijika}

\bksubsection{11}{4分の1波長の線でインピーダンスを変える}
\label{ph:電-13:sec11}
損失のない線の長さを、ちょうど4分の1波長（$\beta l=\pi/2$）にします。$\cos(\pi/2)=0$、$\sin(\pi/2)=1$ なので、行列は1行め $(0,\ jZ_0)$、2行め $(j/Z_0,\ 0)$ です。先に $Z_{\text{先}}$ をつないだとき、入口から見たインピーダンスは、③の式から
\pheq{Z_{\text{入}}=\frac{0\times Z_{\text{先}}+jZ_0}{\dfrac{j}{Z_0}Z_{\text{先}}+0}=\frac{Z_0^{\,2}}{Z_{\text{先}}}}{Ω ＝ Ω$^2$÷Ω}
\label{ph:電-13:lambda4}です。\textbf{4分の1波長の線は、インピーダンスを「ひっくり返す」}はたらきをします。先を開けると（$Z_{\text{先}}\to\infty$）入口はショートに、先をショートすると入口は開いたのと同じになります。

\begin{bkfigure}
\begin{tikzpicture}[x=1mm, y=1mm]
  \draw[bkInk, line width=2.0pt] (0,8) -- (30,8);
  \draw[bkGray!70, line width=1.0pt] (0,8) -- (30,8);
  \node[font=\sffamily\footnotesize, anchor=south] at (15,10) {$Z_1$ の線};
  \filldraw[fill=bkYellow!40, draw=bkInk, line width=0.7pt] (30,5) rectangle (62,11);
  \node[font=\sffamily\footnotesize, anchor=south] at (46,11.5) {$Z_{\mathrm m}=\sqrt{Z_1Z_2}$ の線};
  \draw[bkInk, line width=0.4pt, {Stealth[length=1.6mm]}-{Stealth[length=1.6mm]}] (30,1) -- (62,1);
  \node[font=\sffamily\footnotesize, anchor=north] at (46,0.5) {長さ $\lambda/4$};
  \filldraw[fill=bkGreen!15, draw=bkInk, line width=0.7pt, rounded corners=1mm] (62,1) rectangle (80,15);
  \node[font=\sffamily\footnotesize, align=center] at (71,8) {$Z_2$ の\\機器};
  \draw[bkOrange!85!black, line width=1.0pt, -{Stealth[length=2mm]}] (46,20) -- (34,20);
  \node[font=\sffamily\footnotesize, anchor=south, align=center] at (40,20.5) {ここから見ると $Z_{\mathrm m}^{\,2}/Z_2=Z_1$};
  \node[anchor=west, font=\small, text width=40mm, align=left] at (86,8) {%
  波-7 ⑧の整合層と同じ：$Z_{\mathrm m}=\sqrt{Z_1Z_2}$、厚さ（長さ）は4分の1波長。};
\end{tikzpicture}
\bkfigcaption{4分の1波長の線で、$Z_1$ の線と $Z_2$ の機器を整合させる。}\label{ph:電-13:fig-lambda4}
\end{bkfigure}

\noindent $Z_1$ の線と $Z_2$ の機器のあいだに、特性インピーダンス $Z_{\mathrm m}$、長さ4分の1波長の線をはさむと、入口から見たインピーダンスは $Z_{\mathrm m}^{\,2}/Z_2$ です。これが $Z_1$ に等しければ、はね返りはなくなります。
\[Z_{\mathrm m}=\sqrt{Z_1Z_2}\]
で、波-7 ⑧の超音波のプローブの整合層、メガネの反射防止膜とまったく同じ式です（\bkreidairef{11}）。波-7 では「2つの面ではね返る波が打ち消し合う」と考え、ここでは「四端子定数でインピーダンスを計算する」と考えましたが、同じ答えになりました。テレビのアンテナと機器のあいだの部品や、携帯電話の基地局のアンテナの中でも、この4分の1波長の線が使われています。

\begin{phsaki}{電線の波から、マイクロ波工学・電力系統へ}
この章では、1本の線の上の波を考えました。スマホやレーダーの回路のように、周波数がさらに高い（数 GHz 以上）世界では、線の上の波の反射と整合を、円の図（スミスチャート）で読みながら設計する「マイクロ波工学」が使われます。パソコンの中の高速の信号（1秒に数十億回変わる）の配線の設計でも、基板の線の $Z_0$ と、はしの抵抗によるはね返り止めが欠かせません。電力の世界では、雷が送電線に落ちたときに線の上を走る大きな電圧の波（サージ）を、この章と同じ電信方程式で計算し、変電所の機器を守る設計をしています（電-16）。また、光ファイバーや、電-15 ⑩の導波管の中を進む電磁波も、「特性インピーダンス」と「反射」の考えで扱えます。電線の上の波は、マクスウェル方程式（電-14）からよく計算でき、この章の中身には、わかっていない大きななぞはほとんど残っていません。そのかわり、それをどう使いこなして、速く・遠くへ・むだなく信号や電力を送るかが、今も工学の大きな課題です。
\end{phsaki}

\begin{senmonsho}[電力工学の専門書の書き方]
\begin{bkhonbun}
線路の単位長あたりの直列インピーダンスを $\dot Z=r+j\omega L$，並列アドミタンスを $\dot Y=g+j\omega C$ とすると，伝搬定数 $\dot\gamma=\sqrt{\dot Z\dot Y}=\alpha+j\beta$，特性インピーダンス $\dot Z_0=\sqrt{\dot Z/\dot Y}$ である．送電端電圧・電流 $\dot E_s$，$\dot I_s$ と受電端電圧・電流 $\dot E_r$，$\dot I_r$ の間には
\[\dot E_s=\dot E_r\cosh\dot\gamma l+\dot Z_0\dot I_r\sinh\dot\gamma l,\qquad \dot I_s=\frac{\dot E_r}{\dot Z_0}\sinh\dot\gamma l+\dot I_r\cosh\dot\gamma l\]
の関係がある．
\end{bkhonbun}
\bkyaku{電力の本は、複素数（フェーザ）であることを、文字の上の点「$\dot{\ }$」で示すことが多い（「イー・ドット」と読む）。$r$・$L$・$g$・$C$ は1 m（または1 km）あたりの量で、この章の $R'$・$L'$・$G'$・$C'$。電圧の文字は $E$、送電端（sending end）に添字 s、受電端（receiving end）に r。式は⑨の行列を2行に分けて書いたもの。}
\end{senmonsho}

\begin{taisetsu}
\def\tsl#1{\makebox[7.6\zw][l]{\sffamily\bfseries\small\color{\bkTT}#1}}
\tsl{長い線}長さが波長 $\lambda=c/f$ の約10分の1をこえたら、波として扱う\\[3pt]
\tsl{四端子定数}$\begin{pmatrix}V_1\\I_1\end{pmatrix}=\mathbf F\begin{pmatrix}V_2\\I_2\end{pmatrix}$、つなぐとかけ算、$AD-BC=1$\\[3pt]
\tsl{短い線}$\Delta V\approx I(R\cos\phi+X\sin\phi)$、短絡電流 $I_{\mathrm s}=I_{\mathrm n}\times\dfrac{100}{\%Z}$\\[3pt]
\tsl{電信方程式}$\dfrac{\partial v}{\partial x}=-R'i-L'\dfrac{\partial i}{\partial t}$、$\dfrac{\partial i}{\partial x}=-G'v-C'\dfrac{\partial v}{\partial t}$\\[3pt]
\tsl{損失のない線}$c=\dfrac{1}{\sqrt{L'C'}}$、$Z_0=\sqrt{\dfrac{L'}{C'}}$。同軸は $c=\dfrac{c_0}{\sqrt{\varepsilon_{\mathrm r}}}$、$Z_0\approx\dfrac{60}{\sqrt{\varepsilon_{\mathrm r}}}\ln\dfrac ba$\\[3pt]
\tsl{反射}$r_V=\dfrac{Z_{\text{先}}-Z_0}{Z_{\text{先}}+Z_0}$（波-7 の $r_{\text{押}}$）、はね返る割合 $r_V^{\,2}$\\[3pt]
\tsl{長さ $l$ の線}$\mathbf F=\begin{pmatrix}\cosh\gamma l&Z_0\sinh\gamma l\\\frac{1}{Z_0}\sinh\gamma l&\cosh\gamma l\end{pmatrix}$、$\lambda/4$ の線は $Z_{\text{入}}=\dfrac{Z_0^{\,2}}{Z_{\text{先}}}$
\end{taisetsu}

\begin{kotoba}
\small
\bkword{特性インピーダンス}{線の上を一方向に進む波の、電圧÷電流。線の形と材料で決まり、長さによらない。}
\bkword{二端子対回路}{入口2本・出口2本の線をもつ箱。四端子定数（F 行列）で表す。}
\bkword{縦続接続}{前の箱の出口を、次の箱の入口につなぐこと。行列のかけ算になる。}
\bkword{百分率インピーダンス}{定格の電流で中で下がる電圧が、定格の電圧の何\,\% かで表したインピーダンス。}
\bkword{分布定数回路}{抵抗・インダクタンス・静電容量が、線全体に広がっている回路。}
\bkword{電信方程式}{長い線の上の電圧と電流が満たす、2つの偏微分方程式の組。}
\bkword{整合}{先のインピーダンスを線の特性インピーダンスにそろえて、はね返りをなくすこと。}
\bkword{定在波比}{線の上の電圧の振幅の、いちばん大きい所と小さい所の比。整合していれば1。}
\end{kotoba}

\bkkigocards{\bkkigocard{$Z_0$}{ゼット・ゼロ}{特性インピーダンス〔Ω〕。線の上を進む波の電圧÷電流。電-15 の波のインピーダンスは $\eta$}{characteristic impedance}{電-13}}
{\bkkigocard{$L'$}{エル・ダッシュ}{1 m あたりのインダクタンス〔H/m〕。$R'$・$G'$・$C'$ も「1 m あたり」。$C'$ は四端子定数の $C$ とは別}{inductance per unit length}{電-13}}
\bkkigocards{\bkkigocard{$\gamma$}{ガンマ}{伝搬定数〔1/m〕。$\gamma=\alpha+j\beta$（減衰定数・位相定数）。第2巻の比熱比 $\gamma$ とは別}{propagation constant}{電-13}}
{\bkkigocard{$r_V$}{アール・ブイ}{電圧の反射係数。波-7 の $r_{\text{押}}$ と同じ式。電流の反射係数 $r_I=-r_V$ は波-7 の $r$ と同じ式}{voltage reflection coefficient}{電-13}}

%% ======================== 例題 ========================
\Needspace{10\baselineskip}\bksection{reidai}{例題でたしかめよう}
\label{ph:電-13:ex1}
\begin{phreidai}[同軸ケーブルの $L'$・$C'$・$Z_0$・速さ]
内導体の半径 $a=\qty{0.45}{mm}$、外導体の内側の半径 $b=\qty{1.6}{mm}$、あいだをポリエチレン（比誘電率2.3）で満たした同軸ケーブルがあります（電-3 の例題7 と同じ寸法）。1 m あたりのインダクタンス $L'$ と静電容量 $C'$、特性インピーダンス $Z_0$、信号の速さ $c$ を求めましょう。
\phbamen
細い同軸ケーブル（無線機や測定器に使う形）の値を、形と材料から求める場面です。
\phmushi
導体の抵抗と絶縁体のもれ（損失のない線）。内導体の中の磁場（周波数が高く、電流は表面を流れる）。
\phzu
図\ref{ph:電-13:fig-katachi}の同軸の断面。$a=\qty{0.45e-3}{m}$、$b=\qty{1.6e-3}{m}$、$\varepsilon_{\mathrm r}=2.3$。
\phriyu
同軸の $C'$（電-3 ⑩）と $L'$（電-8 ⑥）を求め、損失のない線の式 $Z_0=\sqrt{L'/C'}$、$c=1/\sqrt{L'C'}$（⑦）に入れます。
\phtateru
\[L'=\frac{\mu_0}{2\pi}\ln\frac ba,\qquad C'=\frac{2\pi\varepsilon_{\mathrm r}\varepsilon_0}{\ln(b/a)},\qquad Z_0=\sqrt{\frac{L'}{C'}},\qquad c=\frac{1}{\sqrt{L'C'}}\]
\phtoku
$b/a=1.6/0.45\approx3.56$ なので、$\ln(b/a)\approx1.27$。$\mu_0/(2\pi)=2\times10^{-7}$ なので、
\[L'=2\times10^{-7}\times1.27\approx\qty{2.5e-7}{H/m}=\qty{0.25}{\micro H/m}\]
\[C'=\frac{2\pi\times2.3\times8.85\times10^{-12}}{1.27}\approx\qty{1.01e-10}{F/m}=\qty{101}{pF/m}\]
\[Z_0=\sqrt{\frac{2.54\times10^{-7}}{1.01\times10^{-10}}}\approx\qty{50}{\ohm}\]
\[c=\frac{1}{\sqrt{2.54\times10^{-7}\times1.01\times10^{-10}}}\approx\qty{2.0e8}{m/s}\]
\phtashikame
単位：$\sqrt{\text{H/F}}$＝Ω、$1/\sqrt{\text{H·F}}$＝1/s に m をかけて m/s（⑦）。\\
大きさ：⑦の $c=c_0/\sqrt{\varepsilon_{\mathrm r}}$ と $Z_0\approx(60/\sqrt{\varepsilon_{\mathrm r}})\ln(b/a)$ で計算しても、約\qty{1.98e8}{m/s}と約\qty{50}{\ohm}で同じです。光の速さの0.66倍で、細い\qty{50}{\ohm}のケーブルの表の値（速度係数66\,\%）と合います。\\
極端な場合：$b\to a$（すきまがなくなる）なら $\ln(b/a)\to0$ で $Z_0\to0$、ショートと同じです。
\kotae $L'\approx\qty{0.25}{\micro H/m}$、$C'\approx\qty{101}{pF/m}$、$Z_0\approx\qty{50}{\ohm}$、$c\approx\qty{2.0e8}{m/s}$（光の約0.66倍）
\end{phreidai}
\begin{jibun}
例題1と同じ寸法で、あいだを発泡ポリエチレン（比誘電率1.5）にすると、$C'$・$Z_0$・$c$ はいくらになりますか。
\bkjans{$C'\approx\qty{66}{pF/m}$、$Z_0=50\times\sqrt{2.3/1.5}$ で約\qty{62}{\ohm}、$c=c_0/\sqrt{1.5}$ で、約 $\qty{2.4e8}{m/s}$（光の約0.82倍）}
\end{jibun}

\label{ph:電-13:ex2}
\begin{phreidai}[75 Ω の同軸ケーブルの太さの比]
テレビ用の\qty{75}{\ohm}の同軸ケーブルを、発泡ポリエチレンで作ります。発泡ポリエチレンの比誘電率を1.5とします。外導体の内側の半径 $b$ は、内導体の半径 $a$ の何倍にすればよいでしょう。
\phbamen
特性インピーダンスが決まった値になるように、ケーブルの形を設計する場面です（はてな？）。
\phmushi
損失、導体の表面のでこぼこ、外導体の網のすきま。
\phzu
同軸の断面。$Z_0=\qty{75}{\ohm}$、$\varepsilon_{\mathrm r}=1.5$、求めるのは $b/a$。
\phriyu
同軸の $Z_0$ は、太さの比の対数と比誘電率だけで決まる（⑦）ので、その式を $b/a$ について解きます。
\phtateru
\[Z_0=\frac{60}{\sqrt{\varepsilon_{\mathrm r}}}\ln\frac ba\qquad\text{より}\qquad\ln\frac ba=\frac{Z_0\sqrt{\varepsilon_{\mathrm r}}}{60}\]
\phtoku
\[\ln\frac ba=\frac{75\times\sqrt{1.5}}{60}=\frac{75\times1.225}{60}\approx1.53,\qquad\frac ba=e^{1.53}\approx4.6\]
\phtashikame
単位：Ω÷Ω で、比には単位がありません。\\
大きさ：内導体が直径\qty{1}{mm}なら、外導体の内側の直径は約\qty{4.6}{mm}で、「5C-FB」（約\qty{5}{mm}）と同じくらいです。\\
極端な場合：比誘電率が大きい絶縁体では、同じ $Z_0$ にするのに、比をもっと大きくしなければなりません。
\kotae 約4.6倍
\end{phreidai}
\begin{jibun}
（1）ポリエチレン（比誘電率2.3）で\qty{50}{\ohm}のケーブルを作るときの太さの比を求めましょう。（2）あいだが空気（比誘電率1）の\qty{50}{\ohm}の線では、比はいくらですか。
\bkjans{（1）$\ln(b/a)=50\times\sqrt{2.3}\div60\approx1.26$ で、約3.5倍（例題1 の寸法）。（2）$\ln(b/a)=50\div60\approx0.83$ で、約2.3倍}
\end{jibun}

\label{ph:電-13:ex3}
\begin{phreidai}[50 Ω のケーブルに 75 Ω の機器をつなぐ]
特性インピーダンス\qty{50}{\ohm}のケーブルの先に、入力のインピーダンスが\qty{75}{\ohm}の機器をつなぎます。電圧の反射係数、電流の反射係数、はね返る電力の割合、機器にかかる電圧（来た波の何倍か）を求めましょう。
\phbamen
波-7 の例題7 と同じ場面を、電気の線の式で確かめます。
\phmushi
ケーブルの中での弱まり。機器のインピーダンスの周波数による変化（\qty{75}{\ohm}のまま）。
\phzu
$Z_0=\qty{50}{\ohm}$、$Z_{\text{先}}=\qty{75}{\ohm}$。図\ref{ph:電-13:fig-jikan}。
\phriyu
先での電圧と電流のつながりから導いた電圧の反射係数（⑧）を使います。電流の反射係数は $-r_V$、電圧は $1+r_V$ 倍。
\phtateru
\[r_V=\frac{Z_{\text{先}}-Z_0}{Z_{\text{先}}+Z_0},\qquad r_I=-r_V,\qquad\text{はね返る割合}=r_V^{\,2},\qquad\frac{v}{v_{\text{来}}}=1+r_V\]
\phtoku
\[r_V=\frac{75-50}{75+50}=\frac{25}{125}=0.20,\qquad r_I=-0.20\]
\[r_V^{\,2}=0.040,\qquad1+r_V=1.2\]
\phtashikame
単位：Ω÷Ω で単位なし。\\
大きさ：機器に伝わる電力の割合は、1から0.040をひいて0.96です。⑧の式で計算しても、$1.2^2\times50\div75$ で0.96になり、合います。波-7 の例題7（$r_V=0.20$、4\,\%）と同じ答えです。\\
極端な場合：$Z_{\text{先}}=\qty{50}{\ohm}$ なら $r_V=0$ で、はね返りません。
\kotae $r_V=0.20$、$r_I=-0.20$、はね返る電力4\,\%、機器の電圧は来た波の1.2倍
\end{phreidai}
\begin{jibun}
\qty{75}{\ohm}のケーブルの先に、\qty{50}{\ohm}の機器をつなぐと、電圧の反射係数、はね返る電力の割合、機器にかかる電圧（来た波の何倍か）はいくらですか。
\bkjans{$r_V=(50-75)\div(50+75)=-0.20$、はね返る電力4\,\%、機器の電圧は $1-0.20=0.80$ 倍}
\end{jibun}

\label{ph:電-13:ex4}
\begin{phreidai}[300 Ω のアンテナ線を 75 Ω の入口につなぐ]
昔のテレビのアンテナ線（平行2線、\qty{300}{\ohm}）を、\qty{75}{\ohm}のテレビの入口に直接つなぎます。電圧の反射係数と、はね返る電力の割合を求めましょう。
\phbamen
波-7 のじぶんでとこう（\qty{300}{\ohm}と\qty{75}{\ohm}）と同じ場面です。
\phmushi
線の弱まり。つなぎ目の部品の大きさ。
\phzu
$Z_0=\qty{300}{\ohm}$、$Z_{\text{先}}=\qty{75}{\ohm}$。
\phriyu
⑧の電圧の反射係数を使います。
\phtateru
\[r_V=\frac{Z_{\text{先}}-Z_0}{Z_{\text{先}}+Z_0},\qquad\text{はね返る割合}=r_V^{\,2}\]
\phtoku
\[r_V=\frac{75-300}{75+300}=\frac{-225}{375}=-0.60,\qquad r_V^{\,2}=0.36\]
\phtashikame
大きさ：$Z$ の比が4倍もあるので、電力の36\,\% がはね返り、テレビに入るのは64\,\% だけです。$r_V<0$ なので、はね返る電圧の波は向きが逆になります（ショートに近い側）。\\
極端な場合：比が1に近づくほど、$r_V$ は0に近づきます。
\kotae $r_V=-0.60$、はね返る電力36\,\%
\end{phreidai}
\begin{jibun}
\qty{50}{\ohm}のケーブルの先に、\qty{50}{\ohm}の機器を2台、並列につなぎました。電圧の反射係数と、はね返る電力の割合を求めましょう。
\bkjans{並列で\qty{25}{\ohm}。$r_V=(25-50)\div(25+50)=-\tfrac13\approx-0.33$、はね返る電力 $\tfrac19\approx11\,\%$}
\end{jibun}

\label{ph:電-13:ex5}
\begin{phreidai}[パルスの往復の時間から、先のようすを知る]
例題1の同軸ケーブル（光の速さの0.66倍）が\qty{10}{m}あります。先に何もつないでいないとき、入口で出したパルスが、先ではね返って戻るまでの時間を求めましょう。戻ってくるパルスの向きは、出したパルスと同じですか、逆ですか。
\phbamen
ケーブルの長さや切れた所を、パルスの往復の時間で調べる方法（\textbf{TDR}：時間領域反射法）の場面です。
\phmushi
ケーブルの中での弱まり。パルスの形のくずれ。
\phzu
図\ref{ph:電-13:fig-tdr}。長さ $l=\qty{10}{m}$、往復の距離 $2l$、速さ $c=c_0/\sqrt{2.3}$。
\phriyu
パルスは一定の速さ $c$ で進むので、時間＝距離÷速さ。向きは、先を開けたときの $r_V=+1$（⑧）で決まります。
\phtateru
\[c=\frac{c_0}{\sqrt{\varepsilon_{\mathrm r}}},\qquad t=\frac{2l}{c}\]
\phtoku
\[c=\frac{3.00\times10^8}{\sqrt{2.3}}\approx\qty{1.98e8}{m/s}\]
\[t=\frac{20}{1.98\times10^8}\approx\qty{1.01e-7}{s}=\qty{101}{ns}\]
\phtashikame
単位：m÷(m/s)＝s。\\
大きさ：1 m 進むのに約\qty{5}{ns}。空気の中の光なら\qty{67}{ns}で、ケーブルの中はその1.5倍かかります。\\
極端な場合：先が切れておらず、正しく\qty{50}{\ohm}でとめてあれば、パルスは戻りません。
\kotae 約\qty{101}{ns}後に、同じ向きのパルスが戻る
\end{phreidai}
\begin{jibun}
例題5のケーブル（\qty{1.98e8}{m/s}）を、長いケーブルにかえて調べたら、\qty{150}{ns}後に、出したパルスと逆向きのパルスが戻ってきました。入口から何 m の所に、何が起きていると考えられますか。
\bkjans{距離は $1.98\times10^8\times150\times10^{-9}\div2$ で約\qty{15}{m}。逆向きなので $r_V<0$、ケーブルがつぶれて2つの導体が近づいた（ショートに近い）所があると考えられる}
\end{jibun}

\label{ph:電-13:ex6}
\begin{phreidai}[短い線か、長い線か]
（1）\qty{50}{Hz}の交流の、空気の中を張った線の上の波長を求めましょう（線の上の速さは $c_0$）。（2）\qty{1}{GHz}の信号の、速さ\qty{2.0e8}{m/s}の線の上の波長を求めましょう。家の\qty{20}{m}の配線と、パソコンの基板の上の\qty{0.1}{m}の線は、それぞれ短い線ですか、長い線ですか。
\phbamen
線を、これまでの回路（集中定数回路）として扱ってよいかを判断する場面です（①）。
\phmushi
線の上の速さのわずかな違い（目安を知るのが目的）。
\phzu
$c=f\lambda$。長い線の目安は、長さが波長の約10分の1をこえること。
\phriyu
比べるのは、線の長さと波長です（①）。
\phtateru
\[\lambda=\frac{c}{f}\]
\phtoku
（1）\[\lambda=\frac{3.00\times10^8}{50}=\qty{6.0e6}{m}=\qty{6000}{km}\]
（2）\[\lambda=\frac{2.0\times10^8}{1\times10^9}=\qty{0.20}{m}\]
家の配線\qty{20}{m}は\qty{6000}{km}の30万分の1で、短い線です。基板の線\qty{0.1}{m}は\qty{0.20}{m}の半分で、10分の1（\qty{2}{cm}）よりずっと長いので、長い線です。
\phtashikame
単位：(m/s)÷(1/s)＝m。\\
大きさ：テレビやスマホの電波（数百 MHz〜数 GHz）では、手のひらほどの線でも「長い線」になります。
\kotae （1）約\qty{6000}{km}（2）約\qty{0.20}{m}。家の配線は短い線、基板の線は長い線
\end{phreidai}
\begin{jibun}
テレビの電波（\qty{600}{MHz}）が、発泡ポリエチレン（比誘電率1.5）の同軸ケーブルの中を進むときの波長を求めましょう。長さ\qty{5}{m}のケーブルは、何波長分ですか。
\bkjans{$c=3.00\times10^8\div\sqrt{1.5}$ で、約 $\qty{2.45e8}{m/s}$。$\lambda\approx\qty{0.41}{m}$。$5\div0.41$ で約12波長分（長い線）}
\end{jibun}

\label{ph:電-13:ex7}
\begin{phreidai}[直列の箱と並列の箱をつなぐ]
直列\qty{2}{\ohm}の箱のあとに、並列\qty{2}{\ohm}（$Y=\qty{0.5}{S}$）の箱をつなぎます（7-3⑥と同じ箱）。全体の行列を求め、$AD-BC=1$ を確かめましょう。出口に\qty{2}{\ohm}の抵抗をつないだとき、入口から見たインピーダンスはいくらですか。
\phbamen
小さな箱をつないで、全体のはたらきを行列で求める場面です（③）。
\phmushi
導線の抵抗。周波数（抵抗だけなので、直流でも交流でも同じ）。
\phzu
図\ref{ph:電-13:fig-hako}の箱を2つ縦続接続。直列の箱 $Z=2$、並列の箱 $Y=0.5$、$Z_{\text{先}}=2$。
\phriyu
縦続接続では、入口側から順に行列をかけます。入口から見たインピーダンスは、③の $Z_{\text{入}}$ の式で求めます。
\phtateru
\[\mathbf F=\begin{pmatrix}1&2\\0&1\end{pmatrix}\begin{pmatrix}1&0\\0.5&1\end{pmatrix},\qquad Z_{\text{入}}=\frac{AZ_{\text{先}}+B}{CZ_{\text{先}}+D}\]
\phtoku
\[\mathbf F=\begin{pmatrix}1+2\times0.5&0+2\times1\\0+1\times0.5&0+1\times1\end{pmatrix}=\begin{pmatrix}2&2\\0.5&1\end{pmatrix}\]
\[AD-BC=2\times1-2\times0.5=1,\qquad Z_{\text{入}}=\frac{2\times2+2}{0.5\times2+1}=\frac{6}{2}=\qty{3}{\ohm}\]
\phtashikame
回路で確かめます。出口の\qty{2}{\ohm}と並列の\qty{2}{\ohm}の合成は\qty{1}{\ohm}（電-4 ⑦）、それに直列の\qty{2}{\ohm}をたして\qty{3}{\ohm}。行列の答えと同じです。\\
単位：$B$ は Ω、$C$ は S で、$AD-BC$ は単位なし。
\kotae $\mathbf F$ は1行め $(2,\ 2)$、2行め $(0.5,\ 1)$、$AD-BC=1$、$Z_{\text{入}}=\qty{3}{\ohm}$
\end{phreidai}
\begin{jibun}
つなぐ順を逆にして、並列\qty{2}{\ohm}の箱のあとに直列\qty{2}{\ohm}の箱をつなぎます。全体の行列、$AD-BC$、出口に\qty{2}{\ohm}をつないだときの $Z_{\text{入}}$ を求めましょう。
\bkjans{$\mathbf F$ は1行め $(1,\ 2)$、2行め $(0.5,\ 2)$。$AD-BC=1\times2-2\times0.5=1$。$Z_{\text{入}}=(1\times2+2)\div(0.5\times2+2)=\tfrac43\approx\qty{1.3}{\ohm}$（\qty{4}{\ohm}と\qty{2}{\ohm}の並列と同じ）}
\end{jibun}

\label{ph:電-13:ex8}
\begin{phreidai}[配電線の電圧降下]
\qty{6.6}{kV}の三相の配電線で、1相分の線の抵抗 $R$ を\qty{0.5}{\ohm}、リアクタンス $X$ を\qty{0.4}{\ohm}とします。電流\qty{100}{A}、力率0.8（電流がおくれる）の負荷に電気を送るとき、1相分の電圧降下と、1相の電圧 $6600/\sqrt3\approx\qty{3810}{V}$ に対する割合を求めましょう。
\phbamen
変電所から工場まで、数 km の配電線で電気を送る場面です（④）。
\phmushi
線のあいだの静電容量（短い線）。三相のつり合いのずれ（1相分で考える：電-11 ⑦）。
\phzu
図\ref{ph:電-13:fig-mijikai}。$R=0.5$、$X=0.4$、$I=100$、$\cos\phi=0.8$。
\phriyu
短い線は直列の箱1つ。電圧降下は $I(R\cos\phi+X\sin\phi)$（④）。
\phtateru
\[\sin\phi=\sqrt{1-\cos^2\phi},\qquad\Delta V\approx I(R\cos\phi+X\sin\phi)\]
\phtoku
\[\sin\phi=\sqrt{1-0.64}=0.6\]
\[\Delta V\approx100\times(0.5\times0.8+0.4\times0.6)=100\times0.64=\qty{64}{V}\]
\[\frac{64}{3810}\approx0.017=1.7\,\%\]
\phtashikame
単位：A×Ω＝V。\\
大きさ：複素数で $V_{\mathrm s}=V_{\mathrm r}+(R+jX)I$ を計算して大きさの差をとっても、64.0 V で、近似の式とほとんど同じです。線間の電圧（2本の線のあいだ）では、$\sqrt3$ 倍の約\qty{111}{V}下がります。\\
極端な場合：$I=0$（負荷なし）なら、電圧は下がりません。
\kotae 約\qty{64}{V}、約1.7\,\%
\end{phreidai}
\begin{jibun}
例題8の配電線で、（1）力率1.0・電流\qty{100}{A}、（2）力率0.8・電流\qty{200}{A} のときの1相分の電圧降下を求めましょう。
\bkjans{（1）$100\times0.5=\qty{50}{V}$。（2）$200\times0.64=\qty{128}{V}$（約3.4\,\%）}
\end{jibun}

\label{ph:電-13:ex9}
\begin{phreidai}[変圧器の出口でショートしたら]
電柱の上の変圧器（単相）の定格を、容量\qty{30}{kVA}、出口の電圧\qty{210}{V}、百分率インピーダンス5\,\% とします。出口のすぐ先でショートしたときに流れる短絡電流を求めましょう。
\phbamen
家や店に電気を送る変圧器で事故が起きたときの電流を見積もり、ブレーカーを選ぶ場面です（⑤）。
\phmushi
変圧器より電源側（配電線・変電所）のインピーダンス（電源はとても強いとする）。ショートした所の抵抗。
\phzu
定格容量 $S_{\mathrm n}=V_{\mathrm n}I_{\mathrm n}=\qty{30}{kVA}$、$V_{\mathrm n}=\qty{210}{V}$、$\%Z=5$。
\phriyu
短絡電流は $I_{\mathrm n}\times100/\%Z$（⑤）。まず定格電流を求めます。
\phtateru
\[I_{\mathrm n}=\frac{S_{\mathrm n}}{V_{\mathrm n}},\qquad I_{\mathrm s}=I_{\mathrm n}\times\frac{100}{\%Z}\]
\phtoku
\[I_{\mathrm n}=\frac{30000}{210}\approx\qty{143}{A},\qquad I_{\mathrm s}=143\times\frac{100}{5}\approx\qty{2860}{A}\approx\qty{2.9}{kA}\]
\phtashikame
変圧器の中のインピーダンスで確かめます。$Z=0.05\times210\div143\approx\qty{0.0735}{\ohm}$ で、$210\div0.0735\approx\qty{2860}{A}$ と同じです。\\
大きさ：定格の20倍。電源側のインピーダンスを入れると、実際はこれより少し小さくなります（この答えは上限の見積もり）。
\kotae 約\qty{2.9}{kA}（定格の20倍）
\end{phreidai}
\begin{jibun}
例題9の変圧器で、百分率インピーダンスが2.5\,\% なら、短絡電流は定格の何倍、何 kA ですか。
\bkjans{$100\div2.5=40$ 倍、$143\times40$ で約\qty{5.7}{kA}}
\end{jibun}

\label{ph:電-13:ex10}
\begin{phreidai}[長い送電線の先の電圧（フェランチ効果）]
長さ\qty{300}{km}の\qty{50}{Hz}の架空送電線で、負荷を切り、先に何もつながないとき、受電端の電圧は送電端の電圧の何倍になりますか。線の損失は無視し、線の上の速さは\qty{3.00e8}{m/s}とします。
\phbamen
夜など、使う電気が少ないときの長い送電線（⑩）。
\phmushi
線の抵抗ともれ（損失のない線）。大地や鉄塔の影響。
\phzu
損失のない線の行列（⑨）。$I_2=0$。$\beta=2\pi/\lambda$、$\lambda=\qty{6000}{km}$。
\phriyu
長い線では、波長に比べた長さが効くので、線の行列 $\cos\beta l$ を使います（⑩）。
\phtateru
\[V_1=\cos\beta l\times V_2,\qquad\frac{V_2}{V_1}=\frac{1}{\cos\beta l},\qquad\beta l=2\pi\,\frac{l}{\lambda}\]
\phtoku
\[\beta l=2\pi\times\frac{300}{6000}=\frac{\pi}{10}\approx0.314\ \text{rad}\ (18°),\qquad\frac{V_2}{V_1}=\frac{1}{\cos18°}=\frac{1}{0.951}\approx1.05\]
\phtashikame
単位：$\beta l$ は rad で単位なし。\\
大きさ：5\,\% 上がります。\qty{500}{kV}の線なら約\qty{26}{kV}も高くなり、機器の絶縁に効きます。\\
極端な場合：$l\to0$ なら $\cos0=1$ で、短い線と同じく電圧は変わりません。
\kotae 約1.05倍
\end{phreidai}
\begin{jibun}
（1）長さ\qty{600}{km}の線では、受電端の電圧は何倍ですか。（2）\qty{50}{Hz}で、受電端の電圧が限りなく大きくなる長さは何 km ですか。
\bkjans{（1）$\beta l=\pi/5$（36°）、$1\div\cos36°$ で約1.24倍。（2）4分の1波長の\qty{1500}{km}（$\cos\beta l=0$）}
\end{jibun}

\label{ph:電-13:ex11}
\begin{phreidai}[4分の1波長の線で整合させる]
\qty{50}{\ohm}のケーブルと、入力のインピーダンスが\qty{100}{\ohm}のアンテナを、周波数\qty{100}{MHz}で整合させたい。あいだに入れる4分の1波長の線の特性インピーダンスと長さを求めましょう。線の上の速さは\qty{2.0e8}{m/s}とします。
\phbamen
アンテナとケーブルのあいだで、はね返りをなくす部品を作る場面です（⑪）。
\phmushi
線の損失。周波数のずれ（\qty{100}{MHz}ちょうどで考える）。
\phzu
図\ref{ph:電-13:fig-lambda4}。$Z_1=50$、$Z_2=100$、$f=\qty{100}{MHz}$、$c=\qty{2.0e8}{m/s}$。
\phriyu
4分の1波長の線は $Z_{\text{入}}=Z_{\mathrm m}^{\,2}/Z_2$ にするので、$Z_{\mathrm m}=\sqrt{Z_1Z_2}$ でそろう（⑪）。
\phtateru
\[Z_{\mathrm m}=\sqrt{Z_1Z_2},\qquad\text{長さ}=\frac{\lambda}{4}=\frac{c}{4f}\]
\phtoku
\[Z_{\mathrm m}=\sqrt{50\times100}=\sqrt{5000}\approx\qty{70.7}{\ohm},\qquad\frac{c}{4f}=\frac{2.0\times10^8}{4\times1.0\times10^8}=\qty{0.50}{m}\]
\phtashikame
入口から見ると $70.7^2/100\approx\qty{50}{\ohm}$ で、ケーブルとそろいます。\\
大きさ：\qty{70.7}{\ohm}のケーブルは特別ですが、\qty{75}{\ohm}のケーブルで代わりにすると、入口は $75^2/100\approx\qty{56}{\ohm}$ で、$r_V\approx0.06$、はね返りは0.3\,\% ほどです。
\kotae 約\qty{70.7}{\ohm}、長さ\qty{0.50}{m}
\end{phreidai}
\begin{jibun}
（1）\qty{50}{\ohm}のケーブルと\qty{200}{\ohm}の機器をつなぐ4分の1波長の線の特性インピーダンスを求めましょう。（2）4分の1波長の線の先をショートすると、入口から見たインピーダンスはどうなりますか。
\bkjans{（1）$\sqrt{50\times200}=\qty{100}{\ohm}$。（2）$Z_0^{\,2}/0$ で限りなく大きくなり、入口は何もつないでいない（開いた）のと同じになる}
\end{jibun}

\label{ph:電-13:ex12}
\begin{phreidai}[もれのある長い線（7-11 の線）]
1 km あたりの抵抗 $R'=\qty{1}{\ohm/km}$、もれ $G'=\qty{1}{S/km}$、長さ\qty{1}{km}の線（7-11⑥と同じ数）に、直流の\qty{1}{V}をかけます。先に何もつながないとき、先の電圧と、入口に流れこむ電流を求めましょう。
\phbamen
もれのある線を、⑨の行列で計算する場面です。数学の 7-11⑥では、箱の数を増やして $\cosh1$・$\sinh1$ に近づくことを見ました。
\phmushi
インダクタンスと静電容量（直流なので効かない：$\omega=0$）。
\phzu
⑨の行列。$Z'=R'=\qty{1}{\ohm/km}$、$Y'=G'=\qty{1}{S/km}$、$l=\qty{1}{km}$、$I_2=0$。
\phriyu
$\gamma=\sqrt{Z'Y'}$、$Z_0=\sqrt{Z'/Y'}$ を求め、長さ $l$ の線の行列に入れます。先を開けるので $I_2=0$。
\phtateru
\[\gamma=\sqrt{R'G'},\qquad Z_0=\sqrt{\frac{R'}{G'}},\qquad V_1=\cosh\gamma l\times V_2,\qquad I_1=\frac{\sinh\gamma l}{Z_0}\times V_2\]
\phtoku
$\gamma=\sqrt{1\times1}=1$〔1/km〕、$Z_0=\qty{1}{\ohm}$、$\gamma l=1$ なので、
\[V_2=\frac{V_1}{\cosh1}=\frac{1}{1.543}\approx\qty{0.648}{V},\qquad I_1=1.175\times0.648\approx\qty{0.762}{A}\]
\phtashikame
行列は1行め $(1.543,\ 1.175)$、2行め $(1.175,\ 1.543)$ で、7-11⑥の表の「限りなく多く」と同じ数です。$AD-BC=1.543^2-1.175^2\approx1$。\\
大きさ：もれがあるので、先へ行くほど電圧が下がります（図\ref{ph:電-13:fig-cosh}）。\\
極端な場合：$G'=0$（もれなし）なら、先を開けると電流は流れず、先の電圧は\qty{1}{V}のままです。
\kotae 先の電圧 約\qty{0.648}{V}、入口の電流 約\qty{0.762}{A}
\end{phreidai}
\begin{jibun}
例題12と同じ線の長さを\qty{2}{km}にすると、先の電圧と入口の電流はいくらですか（$\cosh2\approx3.762$、$\sinh2\approx3.627$）。
\bkjans{$V_2=1\div3.762\approx\qty{0.27}{V}$、$I_1=3.627\times0.266\approx\qty{0.96}{A}$}
\end{jibun}

\begin{bkfigure}
\begin{tikzpicture}
\begin{axis}[bkplot, width=96mm, height=32mm, xmin=0, xmax=1.05, ymin=0, ymax=1.12, xtick={0,0.25,0.5,0.75,1}, ytick={0.648,1}, yticklabels={0.648,1.0}, xlabel={入口からの位置 $x$（km）}, ylabel={電圧（V）}, xticklabel style={/pgf/number format/fixed}]
  \addplot[bkcurve, domain=0:1, samples=60] {cosh(1-x)/cosh(1)};
  \fill[bkRed] (axis cs:0,1) circle (1.8pt) (axis cs:1,0.648) circle (1.8pt);
  \draw[bkGray, line width=0.4pt, dashed] (axis cs:0,1) -- (axis cs:1.05,1);
\end{axis}
\end{tikzpicture}
\bkfigcaption{例題12の線の上の電圧 $V(x)=\cosh(1-x)/\cosh1$〔V〕。もれのために先へ行くほど下がる。先から電流が出ていかないので、先の近くでは傾きが0になる。破線は、もれがないときの\qty{1}{V}。}\label{ph:電-13:fig-cosh}
\end{bkfigure}

%% ======================== まちがえやすいところ ========================
\Needspace{11\baselineskip}\bksection{tsumazuki}{まちがえやすいところをチェック}
\begin{tsumazuki}{50 Ω のケーブルの抵抗？}
\bkng{\qty{50}{\ohm}のケーブルの中の線を、テスターではかると\qty{50}{\ohm}ある。}
\bkok{テスターではかる内導体の抵抗は、長さ\qty{10}{m}でも1 Ω にもならない。\qty{50}{\ohm}は、進む波の電圧÷電流で、長さによらない。}
特性インピーダンス $\sqrt{L'/C'}$ は、熱になって電気をむだにする抵抗とは違う量です（⑦）。
\end{tsumazuki}

\begin{tsumazuki}{先を開けたケーブルに、電流は流れない？}
\bkng{先に何もつながないので、入口からケーブルに電流は流れこまない。}
\bkok{電圧をかけた直後は、ケーブルが $Z_0$ の抵抗のように見えて、電流が流れこむ。先ではね返った波が戻ってきてはじめて、電流が止まる。}
はね返りが戻るまで、入口は先のようすを知りません（⑥の数値実験）。
\end{tsumazuki}

\begin{tsumazuki}{反射係数0.2 なら、20\,\% がはね返る？}
\bkng{$r_V=0.20$ なので、電力の20\,\% がはね返る。}
\bkok{電力の割合は $r_V^{\,2}=0.04$ で、4\,\%。}
$r_V$ は電圧の比、$r_V^{\,2}$ は電力の比です（⑧。波-7 と同じ）。
\end{tsumazuki}

\begin{tsumazuki}{電子が光の3分の2の速さで走る？}
\bkng{ケーブルの中の電子は、\qty{2.0e8}{m/s}で走っている。}
\bkok{電子はとてもゆっくり動くだけで、速く進むのは電圧と電流の形（波）。}
なわの1つ1つの部分が先へ動かないのと同じです（①）。
\end{tsumazuki}

\begin{tsumazuki}{四端子定数の $C$ は静電容量？}
\bkng{四端子定数の $C$ の単位は F（ファラド）。}
\bkok{四端子定数の $C$ は「電流÷電圧」で、単位は S（ジーメンス）。}
同じ文字でも、別の量です。この章では、静電容量は $C'$（1 m あたり）だけで書きました（③）。
\end{tsumazuki}

%% ======================== コラム ========================
\begin{column}{大西洋の海底ケーブルと、電信方程式}
電信方程式の始まりは、海の底に電信のケーブルを通す計画でした。1855年、イギリスのトムソン（のちのケルビン男爵）は、海底ケーブルを「抵抗と静電容量だけの線」と考えて、信号がにぶっておくれる時間は、ケーブルの長さの2乗に比例することを示しました（この式は、熱の伝わり方と同じ形の、波にならない式です：8-13⑥）。1858年に大西洋を横切る最初のケーブルがつながりましたが、信号はとてもおそく、3週間ほどで使えなくなりました。改良したケーブルが1866年に成功し、大西洋の両側が電信でつながりました。

その後、1880年代に、イギリスのヘビサイドが、インダクタンスと絶縁体のもれを入れた、この章の電信方程式の形を整えました。ヘビサイドは、線のインダクタンスを大きくすると、信号がひずまずに遠くまで届くことを示しました（⑨のもっと知りたい）。この考えは、はじめなかなか認められませんでしたが、1900年ごろから、長い電話線にコイルを入れる方法として使われるようになりました。
\end{column}

\begin{column}{電圧が上がる送電線：フェランチ効果の発見}
1890年、ロンドンのテムズ川のほとりのデプトフォードに、若い技術者フェランチが設計した大きな発電所ができました。発電所から街まで、約\qty{11}{km}の地下の同軸のケーブルで、当時としてはとても高い\qty{10}{kV}の交流を送りました。ところが、街の側の電圧が、発電所の側より高くなることがわかり、技術者のあいだで大きな議論になりました。のちに電子管を発明するフレミングらが調べ、ケーブルの大きな静電容量に流れる電流が原因であることがわかりました。これが、⑩の「フェランチ効果」の名前の由来です。
\end{column}

%% ======================== 工学ではこう使う ========================
\Needspace{18\baselineskip}\bksection{phgear}{工学ではこう使う}
\begin{phkougaku}[はね返りで探す・はね返りをなくす]
\textbf{ケーブルの故障を探す}：例題5のように、パルスを送ってはね返りが戻る時間をはかると、長いケーブルのどこが切れたり、つぶれたり、水が入ったりしているかがわかります（TDR）。送電線でも、雷などで事故が起きたとき、線の上を走る電圧の波が変電所に着く時刻の差から、事故の場所を割り出す装置（故障点標定装置）が使われています。

\textbf{パソコンやスマホの基板}：1秒に何十億回も変わる信号を運ぶ基板の線は、特性インピーダンスが決まった値になるように、線のはばと板の厚さを設計し、線の先には $Z_0$ と同じ抵抗をつけて、はね返りを止めます（\textbf{終端}）。はね返りが残ると、信号の0と1を間違えてしまうからです。

\textbf{アンテナとケーブル}：無線機・テレビ・携帯電話の基地局では、アンテナ・ケーブル・機器のインピーダンスをそろえ（整合）、SWR をはかって確かめます。⑪の4分の1波長の線は、そのための部品の1つです。

\textbf{送電線の電圧を保つ}：長い送電線や地中ケーブルでは、④の電圧降下と⑩のフェランチ効果の両方を考えて、変電所でコイルやコンデンサーをつなぎかえ、電圧を決まった範囲に保っています。

\begin{phketa}[線の数の目安]
\phketaline{ケーブルの $Z_0$}{同軸 50・75 Ω、LAN 約100 Ω}
\phketaline{架空送電線の $Z_0$}{約300〜400 Ω}
\phketaline{線の上の速さ}{ケーブル 光の0.6〜0.85倍}
\phketaline{50 Hz の波長}{約\qty{6000}{km}（家の配線は短い線）}
\phboxhead{答えが桁で合っているか}
ケーブルの $Z_0$ が数 Ω や数千 Ω になったら、$L'$ と $C'$ の単位（\si{\micro H}・pF）か、$\ln$ の計算を間違えています。線の上の速さが光の速さをこえたら、間違いです。反射係数の大きさは、抵抗をつないだときはいつも1以下です。
\end{phketa}
\end{phkougaku}
\phstation{電力工学,高周波工学,通信工学,電子回路設計}{特性インピーダンス、整合、反射係数、定在波比（SWR）、四端子定数、縦続接続、伝搬定数、分布定数回路、百分率インピーダンス、フェランチ効果}

%% ======================== 章末問題 ========================
\bkendproblems
\bkproblemgroup{基本}{3}
\bkproblem LAN ケーブルの中の信号の速さを、光の速さの0.70倍とします。長さ\qty{100}{m}の LAN ケーブルを、信号が通りぬけるのにかかる時間を求めましょう（光の速さは\qty{3.00e8}{m/s}）。
\bkpans{$100\div(0.70\times3.00\times10^8)$ で、約 $\qty{4.8e-7}{s}$（\qty{0.48}{\micro s}）}
\bkproblem \qty{75}{\ohm}のケーブルの先に、（1）\qty{100}{\ohm}の機器、（2）何もつながない（開いたまま）、をつなぐと、電圧の反射係数とはね返る電力の割合はいくらですか。
\bkpans{（1）$r_V=25\div175\approx0.14$、はね返る電力 約2.0\,\%。（2）$r_V=+1$、全部（100\,\%）はね返る}
\bkproblem 光の速さの3分の2（\qty{2.0e8}{m/s}）で信号が進むケーブルの入口からパルスを出したら、\qty{0.40}{\micro s}後に、出したパルスと同じ向きのパルスが戻ってきました。入口から何 m の所で、何が起きていると考えられますか。
\bkpans{$2.0\times10^8\times0.40\times10^{-6}\div2=\qty{40}{m}$。同じ向き（$r_V>0$）なので、ケーブルが切れている（開いている）と考えられる}
\bkproblem 無線 LAN の電波（\qty{2.4}{GHz}）の信号が、速さ\qty{2.0e8}{m/s}のケーブルの中を進むときの波長を求めましょう。長さ\qty{1}{m}のケーブルは、短い線ですか、長い線ですか。
\bkpans{$\lambda=2.0\times10^8\div(2.4\times10^9)\approx\qty{0.083}{m}$（\qty{8.3}{cm}）。\qty{1}{m}は約12波長分で、長い線}
\bkproblem 特性インピーダンス\qty{100}{\ohm}の LAN の線の先で、はね返りをなくすには、何 Ω の抵抗をつなげばよいですか。間違えて\qty{120}{\ohm}をつなぐと、電圧の反射係数とはね返る電力の割合はいくらですか。
\bkpans{\qty{100}{\ohm}。\qty{120}{\ohm}なら $r_V=20\div220\approx0.091$、はね返る電力 約0.83\,\%}
\bkproblemgroup{チャレンジ}{5}
\bkproblem 直列\qty{4}{\ohm}の箱のあとに、並列\qty{4}{\ohm}（$Y=\qty{0.25}{S}$）の箱をつなぎます。全体の行列と $AD-BC$ を求め、出口に\qty{4}{\ohm}をつないだときの入口から見たインピーダンスを、行列と回路の2通りで求めましょう（例題7）。
\bkpans{$\mathbf F$ は1行め $(2,\ 4)$、2行め $(0.25,\ 1)$、$AD-BC=2-1=1$。$Z_{\text{入}}=(2\times4+4)\div(0.25\times4+1)=\qty{6}{\ohm}$。回路では、\qty{4}{\ohm}と\qty{4}{\ohm}の並列\qty{2}{\ohm}に直列\qty{4}{\ohm}で\qty{6}{\ohm}}
\bkproblem 発泡ポリエチレン（比誘電率1.5）で\qty{50}{\ohm}の同軸ケーブルを作ります。（1）太さの比 $b/a$ を求めましょう。（2）内導体の半径が\qty{0.50}{mm}なら、外導体の内側の半径はいくらですか。（3）信号の速さを求めましょう（例題1・2）。
\bkpans{（1）$\ln(b/a)=50\times\sqrt{1.5}\div60\approx1.02$、$b/a\approx2.8$。（2）約\qty{1.4}{mm}。（3）$3.00\times10^8\div\sqrt{1.5}\approx\qty{2.4e8}{m/s}$（光の約0.82倍）}
\bkproblem \qty{50}{\ohm}の損失のない線を、\qty{100}{MHz}の8分の1波長の長さにして、先をショートします。入口から見たインピーダンスを⑨の行列と③の式で求め、どんな部品のはたらきか答えましょう。速さを\qty{2.0e8}{m/s}として、長さも求めましょう（例題11）。
\bkpans{$\beta l=\pi/4$、行列は1行め $(\cos\tfrac{\pi}{4},\ j50\sin\tfrac{\pi}{4})$、2行め $(\tfrac{j}{50}\sin\tfrac{\pi}{4},\ \cos\tfrac{\pi}{4})$。$Z_{\text{先}}=0$ なので $Z_{\text{入}}=B\div D=j50\tan\tfrac{\pi}{4}=j\qty{50}{\ohm}$。$j\omega L$ と同じ形なので、コイル（$L=50\div(2\pi\times10^8)\approx\qty{80}{nH}$）のはたらき。長さは $\lambda/8=2.0\times10^8\div10^8\div8=\qty{0.25}{m}$}
\bkproblem ⑨の $V(x)$・$I(x)$ の式が、電信方程式の②（$\dd I/\dd x=-Y'V$）を満たすことを確かめましょう。
\bkpans{$\cosh$ と $\sinh$ を微分すると（中身の微分 $-1$ がつく）、$\dd I/\dd x$ は $-\gamma$ に $(V_2/Z_0)\cosh\gamma(l-x)$ $+I_2\sinh\gamma(l-x)$ をかけたもの。$\gamma/Z_0$ をくくり出すと、残りは $V_2\cosh\gamma(l-x)$ $+Z_0I_2\sinh\gamma(l-x)$ $=V(x)$。$\gamma/Z_0=\sqrt{Z'Y'}\div\sqrt{Z'/Y'}=Y'$ なので、$\dd I/\dd x=-Y'V$}

\begin{furikaeri}
\begin{checklist}
\item ケーブルの上を電気の波が進み、先やつなぎ目ではね返ることを、速さ・往復の時間・波長で説明できる。\bkfurref{ph:電-13:sec1}
\item 特性インピーダンスの違いから、電圧の反射係数と、はね返る電力の割合を求められる。\bkfurref{ph:電-13:sec1}
\item 四端子定数の行列で箱をつなぎ、短い配電線の電圧降下と、変圧器の短絡電流を見積もれる。\bkfurref{ph:電-13:sec3}
\item 小さな箱のくさりから電信方程式を導き、波の速さ $1/\sqrt{L'C'}$ と特性インピーダンス $\sqrt{L'/C'}$ を求められる。\bkfurref{ph:電-13:sec6}
\item 損失のある線・長い送電線・4分の1波長の線を、$\cosh$・$\sinh$ と四端子定数で扱える。\bkfurref{ph:電-13:sec9}
\end{checklist}
\end{furikaeri}
\bkfukidashi{mina}{テレビのケーブルの「75Ω」は、電気の波の手ごたえだったんだね。ひもの波も、音も、電線も、インピーダンスをそろえれば、はね返らない！}
